\documentclass[12pt,a4paper]{report}

% Set lang (support for xelatex and lualatex)
\usepackage{fontspec}
\usepackage{polyglossia}
\setdefaultlanguage{spanish}

% Enhanced font: Computer Modern Unicode
\setmainfont{CMU Serif}
\setsansfont{CMU Sans Serif}
\setmonofont{CMU Typewriter Text}

% Incluye imágenes
\usepackage{graphicx}
\usepackage{tcolorbox} % Es para cuadros de colores
\tcbuselibrary{breakable, skins} % Para poder partir un cuadro
\usepackage{subcaption}

% Math package
\usepackage{amsmath}
\usepackage{amssymb}
\usepackage{amsthm}
% Math graphic
\usepackage{pgfplots}
\usepackage{tikz}
\usepgfplotslibrary{fillbetween}
\usetikzlibrary{positioning, arrows.meta}
\pgfplotsset{compat=1.18}
\usepackage{morefloats}

% Codigo
\usepackage{listings}
\tcbuselibrary{listings} % Es lo que permite usar t colorbox dentro de listings

% Entornos
%=============================
% ENTORNOS PERSONALIZADOS

% Cuadro ejemplo
\newcounter{example}[section]
\renewcommand{\theexample}{\thesection.\arabic{example}}

\newtcolorbox{ejemplo}[1][]{
    enhanced,
    breakable,
    colback=gray!12,
    boxrule=0pt,
    frame hidden,
    borderline west={2pt}{0pt}{black},
    sharp corners,
    left=8pt,
    right=8pt,
    top=6pt,
    bottom=6pt,
    fonttitle=\bfseries,
    title={Ejemplo~\theexample},
    coltitle=black,
    code={\refstepcounter{example}},
    #1
}

% Notas
\newtcolorbox{nota}{
  enhanced,
  breakable,
  colback=gray!6,
  boxrule=0pt,
  frame hidden,
  borderline west={3pt}{0pt}{gray!60},
  sharp corners,
  left=8pt,
  right=8pt,
  top=6pt,
  bottom=6pt,
  fontupper=\small
}

% Ejercicios 
\newtcolorbox{ejercicio}[1][]{
  enhanced,
  breakable,
  colback=white,
  colframe=black!70,      % Color del marco y del título
  colbacktitle=black!70, % Fondo del título
  coltitle=white,        % Texto del título en blanco
  fonttitle=\bfseries\small,
  title={EJERCICIO #1},   % El título es "EJERCICIO" + lo que tú escribas
  detach title,
  % Diseño del encabezado
  before upper={\textsf{\textbf{\textcolor{black!70}{EJERCICIO #1}}}\smallskip\hrule\medskip\par},
  boxrule=0pt,
  frame hidden,
  borderline north={1.5pt}{0pt}{black!70}, % Línea superior gruesa
  borderline south={0.5pt}{0pt}{black!20}, % Línea inferior sutil
  sharp corners,
  left=0pt,   % Alineado a la izquierda sin sangría
  right=0pt,
  top=5pt,
  bottom=5pt,
  fontupper=\small
}


% Teorema
\theoremstyle{definition}
\newtheorem{theorem}{Teorema}%[section]


% ----- CODIGO ----- %

% 1. Definimos los colores para Python
\definecolor{codegreen}{rgb}{0,0.6,0}
\definecolor{codepurple}{rgb}{0.58,0,0.82}
\definecolor{backgray}{rgb}{0.95,0.95,0.96}
\definecolor{pyblue}{rgb}{0.1, 0.4, 0.6}

% 2. Creamos el entorno personalizado llamado "pycode"
\newtcblisting{pycode}[1]{%
    breakable,
    enhanced,
    arc=2pt,
    outer arc=2pt,
    boxrule=0.5pt,
    colback=backgray,
    colframe=pyblue!70!black,
    left=5mm,
    overlay={
        \begin{tcbclipinterior}
            \fill[gray!20] (frame.south west) rectangle ([xshift=5mm]frame.north west);
        \end{tcbclipinterior}
    },
    title=#1, % El título que le pases al entorno
    fonttitle=\bfseries\small\ttfamily,
    listing only,
    listing options={
        language=Python,
        basicstyle=\small\ttfamily,
        keywordstyle=\color{blue}\bfseries,
        stringstyle=\color{codepurple},
        commentstyle=\color{codegreen}\itshape,
        numbers=left,
        numberstyle=\tiny\color{gray},
        stepnumber=1,
        numbersep=8pt,
        showstringspaces=false,
        breaklines=true,
        tabsize=4
    }
}

%=============================

\begin{document}
\begin{titlepage}
    \centering
    \vspace*{\fill}
    {\Huge \ttfamily Analisis de datos en python \par}
    \vspace{2cm}
    {\large \itshape Gino Anci \par}
    {\small\today\par}
    \vspace*{\fill}
\end{titlepage}

\tableofcontents

\chapter{Introducción}

Este documento ha sido diseñado como una guía técnica orientada específicamente a la implementación de código para el análisis de datos. A diferencia de los manuales convencionales, este material no es una introducción a Python. Se asume que el lector posee un dominio fluido de la sintaxis básica, estructuras de control y tipos de datos del lenguaje.

El objetivo aquí no es enseñar a programar, sino transformar la habilidad de programación en una herramienta de exploración científica. Nos enfocaremos en la manipulación eficiente de estructuras multidimensionales, la limpieza técnica de datasets y la traducción de preguntas de negocio en implementaciones algorítmicas utilizando el ecosistema estándar de la industria (Pandas, NumPy, Matplotlib).
\chapter{NumPy}

\section{Introdución}

NumPy (\textit{Numerical Python}) es una libreria de \textit{Python} que se \textit{especializa} al manejo de \textit{arreglos multidimensionales (arrays)} y en la realización de cálculos numéricos de alto rendimiento.

Para el analista, entender NumPy no es solo aprender una librería, sino adoptar una nueva forma de pensar la información ``el pensamiento matricial''.

\subsection{Arreglos multidimensionales: Más allá de las listas}

¿A qué nos referimos exactamente con ``multidimensional''? En términos computacionales, un arreglo multidimensional es una estructura de datos que organiza la información en múltiples ejes o direcciones.

Aunque matemáticamente podemos trabajar con n dimensiones, para facilitar la intuición técnica, solemos categorizarlas en las tres arquitecturas más comunes:
\begin{itemize}
    \item Unidimensional: Conocido como \textit{Vector}.
    \item Bidimensional: Conocido como \textit{Matriz}.
    \item Tridimensional: Conocido como \textit{Tensor}
\end{itemize}


\subsubsection{Una dimensión (Vector)}

\textit{Vector} no es el nombre de tu amigo (o tal vez sí), pero para lo que nos ocupa, es simplemente una lista de elementos dispuestos en una sola dirección. Matemáticamente, un vector se ve así:
\[
V = \left[x_0,x_1,x_2,\dots,x_{n-1}\right]
\]
Ahora bien, a usted lector, probablemente no le interese en lo más mínimo cómo se ve esto en un libro de álgebra lineal; si así fuera, seguramente estaría leyendo álgebra lineal en lugar de este apunte. Para simplificarle la vida, imagine a un vector como una \textit{fila de la escuela}: esa donde el más bajo iba adelante y el más alto atrás. O una fila de asientos en el cine.

Un vector puede representarse como una fila o como una columna de datos. Suponga que tiene una lista de nombres donde cada uno representa a una persona; en este caso, tenemos un \textit{vector de nombres}.

Lo fundamental aquí es que un vector es una estructura unidimensional. No importa si los datos están acostados o parados: solo necesitamos un único \textit{índice} \((i)\) para saber exactamente en qué posición se encuentra un objeto. Si usted sabe que ``Juan'' es el tercero de la fila, no necesita más coordenadas para encontrarlo. En \textit{Python}, ese ``tercer lugar'' es lo único que nos importa para extraer el dato.


\subsubsection{Dos dimensiones (Matriz)}

Una \textit{matriz} es, para fines prácticos, una tabla de filas y columnas. Si usted ha sobrevivido a una hoja de \textit{Excel} o ha intentado jugar un partido de ajedrez, ya conoce esta estructura. Aquí los datos dejan de ser una simple fila india para organizarse en dos direcciones: filas \((i)\) y columnas \((j)\).

Matemáticamente, una matriz se ve de la siguiente manera:
\[
A = \begin{bmatrix}
a_{11} & a_{12} & \cdots & a_{1n} \\
a_{21} & a_{22} & \cdots & a_{2n} \\
\cdots  & \cdots  & \ddots & \vdots\\
a_{m1} & a_{m2} & \cdots & a_{mn}
\end{bmatrix}
\]
Nuevamente, no se asuste con los subíndices; solo están ahí para indicar que ahora necesitamos dos coordenadas para encontrar un dato. Si estuviéramos en un tablero de ajedrez y quisiéramos ubicar a la reina blanca, diríamos que está en la posición \(1,d\). En el mundo de los datos, la lógica es idéntica: primero miramos la fila y luego la columna.

Sin embargo, hay una trampa en la que el programador novato suele caer: en \textit{Python} (y en casi toda la computación moderna), el primer elemento no vive en la posición 1, sino en la 0. Por lo tanto, para acceder al primer dato de nuestra "tabla", utilizaremos la sintaxis [0,0]. Es un pequeño precio a pagar por la eficiencia, aunque a su cerebro le tome un par de días acostumbrarse a no empezar a contar desde el uno.


\subsubsection{Tres dimensiones (Tensor)}

Para no complicarle la vida (ni complicármela yo), imagine un \textit{tensor} simplemente como un cubo o una pila de matrices. Si una matriz es una hoja de cálculo, un tensor es el libro entero de \textit{Excel} donde cada pestaña es una capa de datos. Para ubicar un elemento aquí, ahora necesitamos tres coordenadas: el \textit{piso} (o profundidad), la \textit{fila} y la \textit{columna}.

Básicamente, el mundo físico que habitamos es de tres dimensiones, así que la intuición no debería fallarle.

Tal como usted, no pienso perder el tiempo mostrando cómo se ve un tensor representado matemáticamente; sospecho que ambos tenemos mejores cosas que hacer. Pero si necesita una imagen mental, piense en un \textit{cubo de Rubik}: para identificar un cuadradito específico, usted debe saber en qué cara está (profundidad), en qué fila se encuentra y en qué columna se cruzan.

En el código, esto se traduce en una terna de índices \([z,i,j]\). Donde \(z\) representa el nivel o la capa, e \(i\) y \(j\) son las coordenadas de la ``tabla'' que ya conocimos. Si logra visualizar esto, felicidades, acaba de entender la estructura base que utilizan las redes neuronales modernas para procesar imágenes a color (donde cada "piso" es un canal: Rojo, Verde y Azul).

\subsection{Array}

En programación, un \textit{array} (o "arreglo", si es que usted todavía se resiste al inglés técnico) es una colección de elementos, generalmente números, organizados de una forma específica.

El origen de la palabra, aunque probablemente no le interese, proviene del francés antiguo \textit{arrai}, que significa "poner en orden" o "disponer". Si recupera por un breve instante lo que acaba de leer sobre arreglos multidimensionales, entenderá el punto: los elementos no se tiran al azar dentro de la memoria de la computadora.

Cada dato tiene una posición fija, un \textit{índice} que nos permite encontrarlo rápidamente sin tener que revolver toda la caja. En \textit{NumPy}, esta "disposición" es sagrada. A diferencia de las listas de \textit{Python} convencionales, que son flexibles y un tanto desordenadas, un \textit{array} de \textit{NumPy} es una estructura rígida, eficiente y perfectamente indexada. Si usted sabe el índice, tiene el poder; si no lo sabe, está perdido en la memoria del sistema.

\subsection{¿Por qué usar NumPy?}

Si usted quiere simplificarse la vida y no dejar a su computadora agonizando mientras consume ciclos de CPU innecesarios, le recomiendo encarecidamente que preste atención.

\textit{NumPy} no es solo una opción, es el estándar porque permite:

\begin{itemize}
    \item \textit{Eficiencia a escala:} Trabajar con grandes volúmenes de datos de forma masiva, optimizando el uso de la memoria RAM.
    \item \textit{Velocidad pura:} Los \textit{arrays} son significativamente más rápidos que las listas de \textit{Python}, ya que almacenan los datos de forma contigua en la memoria (algo que a su procesador le encanta).
    \item \textit{Vectorización:} Incluye operaciones que actúan sobre bloques enteros de datos, permitiéndonos jubilar esos bucles \textit{for} lentos y visualmente espantosos.
    \item \textit{Caja de herramientas científica:} Contiene funciones matemáticas avanzadas, herramientas de álgebra lineal, estadística y generadores de números aleatorios.
    \item \textit{Y mucho más:} No voy a nombrar cada beneficio; \textit{NumPy} no me paga por publicitarlo y, siendo honestos, si no fuera software libre, probablemente me cobrarían.
\end{itemize}


\section{Instalación}

Voy a suponer que usted no ha llegado a este apunte sin haber instalado antes otras librerías. En el caso de que sea la primera vez que se enfrenta a la instalación de paquetes, intentaré ser claro (aunque, por favor, espero que ya lo haya hecho antes).

Básicamente, para instalar \textit{NumPy}, al igual que \textit{Pandas} o \textit{Matplotlib}, el procedimiento es el mismo. Dado que asumo que si algún lector saltó directamente a las otras secciones ya dominaba este arte, esta será la única vez que explicaré cómo gestionar paquetes.

Para realizar la instalación, diríjase a la terminal de su editor de código —o si escribe su código directamente en la terminal creyéndose \textit{Linus Torvalds} (en cuyo caso no deberia estar leyendo esto)— y ejecute el siguiente comando:

\begin{verbatim}
pip install numpy
\end{verbatim}

Puede, a su vez, reemplazar \textit{numpy} por \textit{pandas} o \textit{matplotlib} para descargar el resto del ecosistema. Supongo que sabe que es recomendable realizar estas instalaciones dentro de un \textit{entorno virtual}, pero no voy a explicar aquí qué es eso o cómo se configura. Quizás algún día dedique otro apunte a esos menesteres; por ahora, asumo que su entorno está listo para la acción.
\section{Primeros pasos}

Lamentablemente, es hora de dejar la tan hermosa sección de introducción. ¡La voy a extrañar, introducción! Pero como aquí vinimos a procesar datos y no a escribir poesía, es momento de ensuciarnos las manos.

Antes de empezar como locos, hay que importar la librería, porque de lo contrario Python no tendrá idea de qué estamos hablando. Para esto, utilizamos las palabras mágicas: \textit{{\color{red}Import} numpy {\color{red}as} np}. Realmente, la parte de \textit{as np} no es estrictamente obligatoria, pero a la hora de escribir código es mucho más fácil y rápido escribir dos letras que cinco. Si usted decide no usar el alias, prepárese para desperdiciar segundos valiosos de su vida escribiendo \textit{numpy} frente a cada función.

\subsection{Array simple}

Un \textit{array} simple es, esencialmente, el \textit{vector} que vimos en la introducción: una secuencia lineal de datos. Sin embargo, para que \textit{NumPy} haga su magia, este arreglo debe respetar tres reglas de oro:

\begin{itemize}
    \item \textit{Homogeneidad:} Todos los elementos deben ser del mismo tipo. Si un \textit{array} es de números enteros, no puede pretender meter una palabra o un \textit{string} en el medio. Esta rigidez es, precisamente, lo que permite que sea mucho más rápido que las listas convencionales de \textit{Python}.
    \item \textit{Índices:} Cada posición tiene un número de orden que, como ya mencionamos (y espero que no haya olvidado), empieza siempre en \(0\). Por lo tanto, si su \textit{array} tiene \(5\) elementos, el último vive en la posición \(4\).
    \item \textit{Tamaño fijo:} Una vez que se crea el \textit{array} en la memoria, su tamaño queda definido. Aunque existen formas de "alterarlo", lo que realmente sucede tras bambalinas es que se crean copias modificadas y se borran las anteriores. No es tan flexible como una lista, pero es el precio de la eficiencia.
\end{itemize}

\subsubsection{Primer codigo}

Lamentablemente, en \textit{NumPy} nuestro primer código no va a ser el clásico e inútil "Hola mundo"; aquí nuestro bautismo de fuego es el \textit{array simple}. Veamos cómo se traduce a código todo esto que venimos hablando: 

\begin{pycode}{Creación y acceso en un array 1D}
import numpy as np

# Estaturas al azar; el autor no tenía mucha creatividad hoy.
estaturas = np.array([1.75, 1.82, 1.60, 1.91]) 

# Accediendo al primer elemento.
primero = estaturas[0]  # Resultado: 1.75

# Si quisiéramos saber cuántos elementos tiene el arreglo
cantidad = estaturas.size # Resultado: 4
\end{pycode}

Ahora usted me dirá: ``para eso uso una lista de toda la vida''. Pero recuerde, \textit{NumPy} está hecho para cuando tengamos un millón de estaturas y queramos procesarlas en milisegundos. Por cuestiones \textit{didácticas}, no voy a realizar una toma de un millón de datos en este ejemplo; sería una crueldad innecesaria para su memoria RAM. Además, usted lector no se tomaría el trabajo de leer una muestra de ese tamaño, probablemente cerraría este apunte y se iría a ver una serie. Aquí lo que importa es que entienda la lógica, no que vea números infinitos.

\subsubsection{La regla de la homogeneidad}

Como se menciono anteriormente un array tiene todos los espacios contados, si intentaramos meter algo diferente, \textit{NumPy} va a tener que tomar una decisión para que el array siga siendo eficiente.

Usted lector, curioso y quizás algo malvado por naturaleza, seguramente se cuestiona: ¿Qué pasa si intento romper la regla adrede? Si usted crea un \textit{array} mezclando números enteros o flotantes con palabras, \textit{NumPy} transformará a \textit{todos} los elementos al tipo de dato más flexible para que ninguno pierda información. En la jerarquía de datos, esto suele significar que todo terminará convertido en un \textit{string}.

\begin{pycode}{Un Intruso en el array}
    import numpy as np

    edades = np.array([25, 30, 'Gino', 22]) # ¡Gino no es una edad, es mi nombre!

    # Imprimimos edades
    print(edades) # devuelve ['25' '30' 'Gino' '22']
\end{pycode}
Si usted está atento, se habrá dado cuenta de que \textit{NumPy} puso comillas en todos los valores. Básicamente, convirtió a sus pobres números en simples textos.

\textit{NumPy} hace esto para salvarle las papas a su computadora. Recuerde que la librería asigna un espacio fijo en la memoria para cada elemento; si todos son enteros pequeños, reserva pocos \textit{bytes}, pero si aparece un \textit{string}, debe asignar espacio suficiente para el texto más largo. Para mantener la estructura "ordenada" que vimos antes, todos deben medir lo mismo.

Este proceso de conversión automática se denomina \textit{coerción de tipos}, aunque probablemente usted y yo lo olvidemos apenas cerremos este apunte.


\subsubsection{Operaciones vectorizadas}

Esta es la razón por la cual usamos \textit{Numpy} y no otra calculadora o una lista de python.

La \textit{vectorización} es la capacidad de aplicar una operación matemática a un \textit{array} completo, o entre dos \textit{arrays}, sin necesidad de escribir ciclos (\textit{loops}) manuales. Para entenderlo de forma sencilla: imagine que tiene que firmar y poner sus datos en 100 hojas. Puede hacerlo a mano con una lapicera (listas normales), o puede mandar a fabricar un sello con todos sus datos y simplemente presionar el mismo sobre cada hoja. \textit{NumPy} es el sello.

\begin{ejemplo}
    Imagine que usted es un chef y una pareja que va a casarse lo contrata para servir 1.000 platos de fideos; además, a cada uno debe ponerle una montaña de queso rallado.

    Si usara \textit{Python} normal, usted debería ir con un rallador y un bloque de queso al primer plato: ralla, ralla, ralla. Luego se dirige al segundo: ralla, ralla, ralla... y así hasta el plato 1.000. El resultado: para cuando termine, los fideos estarán fríos, usted estará viejo y la boda probablemente habrá terminado hace décadas.

    Pero como usted es un chef inteligente, le pide a la pareja que gaste 60.000 dólares en una máquina industrial con un sistema de ventilación en el techo, cargado con 50 kg de queso. Usted solo aprieta un botón y una ráfaga de aire precisamente calculada deposita el queso en los 1.000 platos al mismo tiempo.

    Lamentablemente, aunque la idea sea buena, probablemente al chef no lo vuelvan a contratar o lo llamen loco. Pero nosotros no somos chefs, somos analistas de datos. \textit{NumPy} es esa máquina industrial: hace el trabajo sucio por nosotros en un solo paso, mientras que de forma manual nos tomaría una eternidad.
\end{ejemplo}

Como aún no hemos avanzado lo suficiente como para escribir un código complejo, veamos cómo se ve esta "magia" en un ejemplo sencillo. Una operación vectorizada se vería así:

\begin{pycode}{Operación vectorizada}
    import numpy as np

    # Creamos la variable platos de fideos, donde el elemento de la lista es el peso del plato en gramos
    platos = np.array([400, 412, 378, 415])

    # Para cada plato rallamos 20 gramos de queso pero con la maquina automatizada sin maquina automatizada seria iterar en el bucle for
    rallar_queso = platos + 20
    print(f'Peso del plato sin queso = ', platos) # [400 412 378 415]

    print(f'Peso del plato con queso = ', rallar_queso) # [420 432 398 435]
\end{pycode}

Veamos ahora que la magia no se limita a sumar un número solitario; también podemos operar entre vectores completos.

\begin{ejemplo}
    La vida de un chef es difícil: no solo tiene que rallar queso, sino que probablemente deba lidiar con cinco platos de distintos tamaños y cinco porciones de salsa. Para obtener el peso total de cada comida, usted no quiere tirar cualquier salsa sobre cualquier pasta. Usted necesita que la salsa [0] vaya al plato [0], la salsa [1] al plato [1], y así sucesivamente.
\end{ejemplo}

Si quisiéramos realizar una operación entre vectores en código, \textit{NumPy} se encarga de la logística de emparejamiento por nosotros:
\begin{pycode}{Operación de vectores}
import numpy as np

# 1. Creamos el vector de "Pasta" (el peso en gramos de 5 platos diferentes)
pasta = np.array([100, 120, 110, 130, 105])

# 2. Creamos el vector de "Salsa" (el peso de la salsa lista para cada plato)
salsa = np.array([50, 60, 55, 65, 52])

# NumPy es inteligente: sabe que debe sumar la posición [0] con la [0], la [1] con la [1], y así hasta el final. No necesitas decirle cómo hacerlo.
plato_completo = pasta + salsa

# --- RESULTADOS ---
print("Pesos de la Pasta:      ", pasta)
print("Pesos de la Salsa:      ", salsa)
print("-----------------------------------------")
print("PESO TOTAL DE CADA PLATO:", plato_completo)
\end{pycode}

Sin embargo, ojalá la vida fuera todo colores y chispitas. Existe una ley inquebrantable en la aritmética de arreglos: para poder realizar una suma, el tamaño del primer vector \textit{debe ser exactamente igual} al tamaño del segundo.

Si usted tiene cinco platos de pasta pero solo cuatro porciones de salsa, \textit{NumPy} no va a "adivinar" qué plato se queda vacío; simplemente lanzará un error (\textit{ValueError}) y detendrá su cocina industrial. En el caso de la multiplicación existen otras reglas de etiqueta matemática, pero eso lo veremos más adelante, cuando sea estrictamente necesario para no saturar su paciencia.
\section{Creación de Arrays}

Hay varias formas de generar un \textit{array}: puede ser a través de listas o tuplas, mediante funciones como \textit{arange}, \textit{linspace}, y más. Sin embargo, usted no está aquí para que yo simplemente se las mencione; usted las quiere ver en acción, así que vayamos a eso.

\subsection{Arrays desde listas o tuplas}

\textit{NumPy} nos permite transformar las estructuras clásicas de \textit{Python} en potentes arreglos utilizando la función que ya hemos estado ``spoileando'' en las secciones anteriores:

\begin{verbatim}
np.array()
\end{verbatim}

Si ya tiene una lista o una tupla y desea dotarla de superpoderes para operar con ella en \textit{NumPy}, lo único que debe hacer es pasarle ese objeto como argumento a la función. Véalo en el siguiente bloque de código:

\begin{pycode}{array desde listas y tuplas}
import numpy as np

# ---- Caso Lista ---- #
# Tenemos una lista
lista = [1,2,3,4,5] # El autor no tenía muchas ganas de pensar nombres hoy.

# La convertimos en array 
array_desde_lista = np.array(lista) # ¡Si no pense el nombre de la lista, menos el del array!

# ---- Caso Tupla ---- #
tupla = (10,20,30,40)
array_desde_tupla = np.array(tupla)

# ---- Sorpresa ---- #
lista_bidimensional = [[1,2,3], [4,5,6]] # ¡Estamos creando una matriz!
array_matriz = np.array(lista_bidimensional)
\end{pycode}

Para su felicidad, usted, mi lector ansioso, ahora puede ver cómo se materializa en código aquel concepto de \textit{matriz} que explicamos en la introducción. Note cómo \textit{NumPy} interpreta automáticamente las listas anidadas como filas de una tabla.

\subsection{Función \textit{np.arange}}

\textit{Python} nativo tiene una función llamada \textit{range}. Pues bien, \textit{np.arange} (que proviene de \textit{array range}, de ahí la "a" adicional) es la versión supercargada de esa función.

Es una herramienta que crea un \textit{array} de una sola dimensión, rellenándolo con números que siguen un orden lógico: una \textit{progresión aritmética}.

Una progresión aritmética es simplemente una serie de números donde la distancia entre un escalón y el siguiente es siempre la misma. Matemáticamente, una progresión con paso tres (3) se vería así:
\[
2,5,8,11,14
\]

Más allá de la matemática, imagine a \textit{np.arange} como una máquina que anota números en una lista, a la cual usted debe darle tres instrucciones: dónde comenzar, dónde terminar y de cuánto en cuánto avanzar. La estructura es:

\begin{verbatim}
np.arange(inicio, fin, paso)
\end{verbatim}

Pero ojo al piojo, el número de parada \textit{nunca} pero \textit{nunca} se incluye en el resultado final. S le pides que pare en \(10\), podria llegar hasta el \(9,9999\dots\) pero nunca tocará el \(10\). Matemáticamente esto se veria como \([inicio,fin)\) donde el ``\([]\)'' es un intervalo cerrado, es decir que los números extremos está dentro, es decir lo incluye, y el ``\(()\)'' es un intervalo abierto, donde uno se puede acercar mucho al valor, pero nunca tocar el borde.

Veámoslo con algunos ejemplos prácticos:
\begin{pycode}{np.arange}
import numpy as np

# Generamos una propersión aritmética solo con el valor donde debe finalizar (automaticamente comienza en 0)
arange_sin_inicio_ni_paso = np.arange(5) # si lo imprimieramos nos devolveria [0, 1, 2, 3, 4]

# Con inicio y fin
arange_sin_paso = np.arange(5, 10) # nos devolveria [5, 6, 7, 8, 9]

# Con todo incluido
arange_completo = np.arange(0,11,2) # devuelve [0, 2, 4, 6, 8, 10]

# Con decimales
arange_decimal = np.arange(0, 1, 0.1) # devuelve decimales del 0.1 al 0.9

\end{pycode}

Usted, lector escéptico, me dirá: "¿Para qué me sirve esto en la vida real?". Pues claramente no le servirá para ir a comprar el pan, pero si usted quisiera realizar un gráfico de ventas mes a mes, podríamos utilizar esta función para generar los números de los meses del 1 al 12 sin complicaciones.

Ahora, si a usted le interesan las finanzas, imagínese que quiere proyectar qué pasa con una inversión frente a una tasa de interés que varía desde el \(1\%\) al \(10\%\), saltando de medio punto en medio punto. En lugar de escribir cada valor a mano como un pasante mal pagado, usted simplemente ejecutaría:

\begin{verbatim}
tasas = np.arange(0.01, 0.105, 0.005)
\end{verbatim}

Fíjese que usamos \(0.105\) como límite final para que \textit{NumPy} incluya el \(0.10\) \((10\%)\) en el resultado. Gracias a la \textit{vectorización} que vimos antes, usted podría calcular el retorno de inversión para todas esas tasas al mismo tiempo con una sola línea de código. La eficiencia no es solo para las máquinas; es para que usted pueda terminar su trabajo antes y dedicarse a cosas más interesantes.



\subsection{Función \textit{np.linspace}}

Si \textit{np.arange} es una máquina que da pasos a tamaños fijos, \textit{np.linspace} o \textit{Linear space} es una máquina que rellena un espacio.

Básicamente, mientras que en \textit{np.arange} generábamos datos definiendo el "salto", a \textit{np.linspace} le decimos cuántos datos queremos dentro de un intervalo. Es decir, \textit{np.linspace} tiene un inicio, un final y una cantidad de elementos que deseamos repartir equitativamente entre ambos puntos. La estructura es:

\begin{verbatim}
np.linspace(inicio, fin, cantidad)
\end{verbatim}

A diferencia de su primo \textit{arange}, aquí hay una regla de cortesía importante, \textit{np.linspace} utiliza intervalos cerrados \([inicio,fin]\), lo que significa que el número de inicio y el de fin sí están incluidos en el resultado. Véalo en código:

\begin{pycode}{np.linspace}
import numpy as np 

# Queremos generar 4 números repartidos entre 0 y 10.
linspace_basico = np.linspace(0, 10, 4) # nos retorna [ 0. , 3.33333333, 6.66666667, 10. ]

\end{pycode}

Usted, lector insistente, me va a preguntar para qué sirve esto. Pues bien, \textit{np.linspace} es el mejor amigo de los gráficos. Si tenemos una función matemática y queremos dibujar una curva suave en lugar de una línea quebrada y espantosa, simplemente le pedimos a \textit{linspace} que genere 100 o 1.000 puntos en el intervalo. A más puntos, más suave será la curva y más profesional se verá su trabajo.


\subsection{Función \textit{np.zeros}}

La función \textit{np.zeros} es como cuando te mudas solo y abres la heladera por primera vez: está totalmente vacía. Para los que vienen del mundo de la oficina, es como abrir una hoja de \textit{Excel} nueva; no tiene datos aún, pero está lista para que usted empiece a rellenarla.

En el análisis de datos, muchas veces no tenemos los datos de inmediato, pero ya sabemos qué estructura van a tener. Esto nos permite reservar un lugar en la memoria de la computadora mediante la función:

\begin{verbatim}
np.zeros()
\end{verbatim}

Esta función crea un \textit{array} donde absolutamente todos los elementos son el número 0. Es la forma más común de "inicializar" un arreglo antes de llenarlo con información real.

\begin{ejemplo}
Imagine que nuestro chef de la boda quiere realizar una encuesta sobre la calidad de los platos. Antes de empezar a preguntar, arma una planilla con 1.000 filas vacías (es decir, con ceros). A medida que cada invitado responde, el chef anota el resultado; por detrás, \textit{NumPy} borra el cero y escribe el nuevo dato.
\end{ejemplo}

Por lo tanto, \textit{np.zeros()} es nuestra herramienta para crear esa "planilla en blanco". Para que funcione, necesita que usted le dé la \textit{forma} (\textit{shape}) dentro de los paréntesis: si quiere un vector, pase un solo número; si quiere una matriz, pase un par de números entre paréntesis (una tupla).

Véalo en código:

\begin{pycode}{np.zeros}
import numpy as np 

# Creamos una lista de un vector con 5 ceros:
vector_vacio = np.zeros(5) # Retorna [0., 0., 0., 0., 0.]

# Creamos una matriz de 2 filas y 3 columnas
matriz_vacia = np.zeros((2, 3)) # OJO, es muy importante el doble parentesis, ya que el interno define la tupla y el de fuera la función.

# La matriz vacia retornaria:
# [[0., 0., 0.],
# [0., 0., 0.]]

# Si quiere enteros:
vector_enteros_vacios = np.zeros(4, dtype=int) # Retorna [0, 0, 0, 0]

\end{pycode}


\subsection{Función \textit{np.ones}}

Si \textit{np.zeros()} creaba una pizarra de ceros, \textit{np.ones()} crea un \textit{array} donde todos los elementos son, lógicamente, el número 1. Esta función es sumamente útil cuando necesitamos inicializar valores que luego servirán como base para multiplicaciones o para crear "máscaras" de datos.

\begin{ejemplo}
    Imagine que tiene un vector de precios y quiere aumentar todos un \(10\%\) Una forma elegante de hacerlo sería crear un vector de unos del mismo tamaño, sumarle \(0.10\) (quedando todos en \(1.10\)) y luego multiplicar ese vector por sus precios originales.
\end{ejemplo}

En cuanto a la sintaxis, funciona exactamente igual que \textit{np.zeros()}, solicitando la "forma" (\textit{shape}) deseada:

Vease en código:

\begin{pycode}{Uso de np.ones}
import numpy as np

vector_unos = np.ones(3) # Retorna [1., 1., 1.]

matriz_unos = np.ones((2, 2))
# Retorna
# [[1., 1.],
# [1., 1.]]

\end{pycode}

\subsection{Función \textit{np.eye}}

Esta función tiene un nombre que cobra sentido en inglés: la palabra \textit{eye} suena exactamente igual que la letra ``I'', que en matemáticas representa a la \textit{Matriz Identidad}. Lamentablemente, en español nos perdemos el juego de palabras, ya que ``ojo'' no tiene ninguna relación fonética con ``identidad''.

Una matriz identidad es una tabla donde todos los números son \(0\), excepto en la \textit{diagonal principal} (aquella que recorre la matriz desde la esquina superior izquierda hasta la inferior derecha). En esa diagonal, mandan los unos. Matemáticamente, se ve así:
\[
\begin{bmatrix}
    {\color{red}1} & 0 & 0 \\
    0 & {\color{red}1} & 0 \\
    0 & 0 & {\color{red}1}
\end{bmatrix}
\]

Más allá de la pizarra, está nuestro código. A diferencia de las funciones anteriores, \textit{np.eye()} suele recibir un solo número, ya que las matrices identidad son, por definición, cuadradas (mismas filas que columnas). Veamos el código:

\begin{pycode}{np.eye}
import numpy as np

identidad = np.eye(3)

"""
Resultado:
[[1., 0., 0.],
[0., 1., 0.],
[0., 0., 1.]]
"""
\end{pycode}

En la aritmética común, el número \(1\) es el "elemento neutro" de la multiplicación: cualquier número multiplicado por \(1\) da el mismo número \((5\times1=5)\).
En el mundo de las matrices (Álgebra Lineal), la Matriz Identidad cumple esa misma función. Si multiplicas cualquier matriz por la Identidad, obtienes la misma matriz original.

¿Para qué nos sirve esto en análisis de datos? Muchas veces multiplicamos matrices para aplicar transformaciones, como rotar un gráfico, aplicar pesos a indicadores financieros o realizar proyecciones complejas donde necesitamos un "punto de partida" neutro.


\section{Propiedades y estructura de los arrays}

\subsection{Dimensiones, forma, tamaño y tipo de datos}

\subsubsection{Dimensiones \textit{ndim}}

Las dimensiones son las que nos permiten entender cómo se organiza la información en el espacio (y en la memoria de su pobre computadora). La propiedad \textit{.ndim} —o \textit{number of dimensions} para el que todavía se pelea con el inglés— nos dice cuántas coordenadas necesita \textit{NumPy} para encontrar un dato específico.

Básicamente, \textit{.ndim} nos informa si nuestro objeto es:

\begin{itemize}
\item \textit{Escalar (0 dimensiones):} Es un simple número suelto, como un 5. Imagine un punto en el vacío: no necesita dirección porque ya está ahí. No se mueva, no busque, solo es.
\item \textit{Vector (1 dimensión):} Es una línea. Necesita una coordenada (el índice) para saber en qué posición de la fila se encuentra el dato.
\item \textit{Matriz (2 dimensiones):} Es un plano o una tabla. Aquí ya necesita dos coordenadas: la fila y la columna. Si solo da una, \textit{NumPy} le entregará una fila entera y se quedará esperando el resto.
\item \textit{Tensor (3 dimensiones):} Es un cubo. Aquí la cosa se pone tridimensional y necesita tres coordenadas: el piso (profundidad), la fila y la columna.
\end{itemize}

Saber el número de dimensiones es fundamental para no romper nada. Si usted intenta sumar un cubo con una línea, \textit{NumPy} le lanzará un error de dimensiones en la cara; usar \textit{.ndim} es su primera línea de defensa para depurar el código antes de que todo explote.

Veamos \textit{.ndim} en acción:

\begin{pycode}{Uso de .ndim en NumPy}
import numpy as np

# Uso de un vector simple 
a = np.array([1, 2, 3])
print(a.ndim)  # Resultado: 1

# Uso de una matriz 
b = np.array([
    [1, 2],
    [3, 4]
    ])
print(b.ndim)  # Resultado: 2

# Uso de un tensor o paquete de matrices
c = np.array([
    [[1], [2]],
    [[3], [4]]
    ])
print(c.ndim)  # Resultado: 3

\end{pycode}

Para aquellos que todavía no terminan de "cazar" las dimensiones visualmente, hay un truco sucio pero infalible: cuente cuántos corchetes ``\textit{[ ]}'' se abren al inicio del array. Si ve tres corchetes seguidos, felicidades, está ante un tensor de tres dimensiones. Si ve uno solo, es un vector. Es una regla simple, elegante y le ahorrará mirar la documentación más de lo estrictamente necesario.



\subsubsection{Forma o \textit{shape} (para los bilingües)}

Si \textit{.ndim} nos decía cuántas dimensiones tiene el objeto, \textit{.shape} nos indica qué tan "grande" es cada una de ellas. Esta propiedad es, básicamente, una radiografía de sus datos.

Técnicamente, \textit{.shape} nos devuelve una \textit{tupla} (esos números entre paréntesis que no se pueden tocar), donde cada valor representa la cantidad de elementos en un eje específico.

Aunque parezca una propiedad inofensiva, la consultará más veces que su saldo bancario, ya que:
\begin{itemize}
\item Nos permite conocer la magnitud de los datos antes de procesarlos.
\item Es el juez que valida operaciones matemáticas: si intenta sumar matrices con \textit{shapes} distintos, la operación fallará miserablemente.
\item Es fundamental para preparar los datos antes de enviarlos a librerías de gráficos como \textit{Matplotlib}.
\end{itemize}

Véalo en esta demostración técnica:

\begin{pycode}{Uso de .shape en NumPy}
import numpy as np

# Un Vector
vector = np.array([10, 20, 30, 40, 50])
print(vector.shape) # Resultado: (5,) es decir 5 elementos en una sola linea. 

# Una Matriz (2D)
matriz = np.array([
    [1, 2, 3],
    [4, 5, 6]
    ])
print(matriz.shape) # Resultado: (2, 3) es decir  2 filas y 3 columnas.

# Un Cubo (3D)
cubo = np.zeros((2, 3, 4)) # Aqui utilizamos .zeros() para no morir escribiendo números
print(cubo.shape) # Resultado: (2, 3, 4) que son 2 matrices, cada una de 3 filas y 4 columnas.

\end{pycode}

Un detalle para su salud mental: note que en el vector el resultado es \textit{(5,)}. Ese espacio vacío después de la coma está ahí para recordarle que solo hay una dimensión. No intente buscar el segundo número, no existe, es solo \textit{NumPy} siendo técnicamente correcto (y un poco molesto).



\subsubsection{Tamaño o \textit{size}}

Este es el último integrante de la "Santísima Trinidad" de la inspección. Mientras que los otros se preocupan por la estética y la estructura, \textit{.size} es un contador pragmático: nos dice cuántos elementos hay en total en el \textit{array}, sin importar cómo estén distribuidos en filas, columnas o dimensiones.

Matemáticamente, \textit{.size} es el producto de todos los números que aparecen en \textit{.shape}. Por ejemplo:

\begin{itemize}
\item En un vector con \textit{.shape} (5,), el \textit{.size} es \(5\).
\item En una matriz con \textit{.shape} (3,4), el \textit{.size} es \(3 \times 4 = 12\).
\item En un tensor con \textit{.shape} (2,2,2), el \textit{.size} es \(2 \times 2 \times 2 = 8\).
\end{itemize}

Usted, lector escéptico, se cuestionará: "¿Realmente necesito saber cuántos datos hay en total?". La respuesta es un rotundo sí, principalmente para no hacer explotar su sistema. El atributo \textit{.size} nos ayuda a tener un control estricto de la memoria y es el juez supremo cuando queremos usar \textit{.reshape()} (lo veremos en la siguiente subsección); usted no puede transformar una tabla de 12 elementos en una de \(5 \times 5\)5×5, porque los datos no aparecen por arte de magia.

En codigo (no muy diferente a los anteriores), esto se ve así:

\begin{pycode}{Diferencia entre .size}
import numpy as np

# Creamos una matriz de 3 filas y 4 columnas

matriz = np.array([
[1, 1, 1, 1],
[2, 2, 2, 2],
[3, 3, 3, 3]
])

print(f"Total de elementos (.size): {matriz.size}") # Resultado: 12
\end{pycode}


\subsubsection{Tamaño, forma y dimensión}

Por si todavía tiene una ensalada de conceptos en la cabeza, aquí tiene el resumen. \textit{.ndim} le dice en cuántas direcciones se mueven los datos; \textit{.shape} le da el mapa exacto de esas direcciones; y \textit{.size} le dice cuántos elementos hay en total.

La diferencia se observa mejor en el código:
\begin{pycode}{Diferencia entre .ndim .shape y .size}
import numpy as np

# Creamos una matriz de 3 filas y 4 columnas

matriz = np.array([
[1, 1, 1, 1],
[2, 2, 2, 2],
[3, 3, 3, 3]
])

print(f"Dimensiones (.ndim): {matriz.ndim}")  # Resultado: 2
print(f"Forma (.shape): {matriz.shape}")      # Resultado: (3, 4)
print(f"Total de elementos (.size): {matriz.size}") # Resultado: 12
\end{pycode}


\subsubsection{El metodo \textit{.reshape()}}

Seguramente usted, lector indignado, saltará a mi yugular diciendo: "¡Esto no es un atributo, es un método! ¿Qué hace en esta sección?". Pues bien, el autor soy yo y considero que explicar \textit{.reshape()} aquí es mucho mejor que hacerlo más adelante, cuando usted probablemente ya haya olvidado qué es el \textit{.shape}.

A diferencia de los atributos que vimos antes (que solo nos daban información), \textit{.reshape()} es una herramienta activa. Nos permite cambiar la estructura de un \textit{array} —sí, lo sé, suena a magia negra— \textit{sin tocar ni uno solo de los datos que tiene dentro}.

\begin{ejemplo}
    Imagine que tiene una tableta de chocolate con 12 cuadraditos dispuestos en una sola fila larga (nuestro chef se tomó un descanso, así que ahora comemos chocolate). \textit{.reshape()} es el acto de romper esa fila y reacomodarla para que ahora tenga 3 filas de 4 cuadraditos. El chocolate es el mismo y el peso es idéntico, pero su forma ha cambiado para que entre en una caja distinta.
\end{ejemplo}

Para utilizar este método, el \textit{.size} cobra más sentido que nunca. La regla de oro es inquebrantable: el número total de elementos debe ser exactamente el mismo antes y después de la transformación. Usted no puede generar datos de la nada ni puede hacer desaparecer cuadraditos de chocolate en el proceso; NumPy no tolera la magia contable.

Véalo en el código:

\begin{pycode}{Uso de .reshape()}
import numpy as np

# Creamos un vector de 12 elementos (del 0 al 11)
vector = np.arange(12) # Forma actual: (12,) Tamaño: 12

# Lo transformamos en una matriz de 3 filas y 4 columnas 
matriz = vector.reshape(3, 4)
print(matriz)
"""
Resultado:
[[ 0,  1,  2,  3],
[ 4,  5,  6,  7],
[ 8,  9, 10, 11]]
"""

print(matriz.shape) # Resultado: (3, 4)
\end{pycode}

\begin{nota}
\textit{El truco del comodín (-1):} Si usted es un poco perezoso (como todo buen programador), NumPy tiene un regalo para usted. Si conoce cuántas columnas quiere, pero no tiene ganas de calcular cuántas filas saldrán, puede usar el \textbf{-1}. Es un mensaje que dice: "Calcula tú este número basándote en el \textit{.size}".
\end{nota}

\begin{pycode}{El comodin -1 en reshape}
import numpy as np 

vector = np.arange(12)
# Quiero 2 filas y no se cuantas columnas
nueva_forma = vector.reshape(2, -1) # Resultado: Matriz de (2, 6)

\end{pycode}


\subsubsection{Tipos de datos \textit{.dtype}}

Si recuerda el inicio de este apunte, mencionamos que un \textit{array} es estrictamente \textit{homogéneo}, todos los elementos deben ser del mismo tipo. Pues bien, con \textit{.dtype} (\textit{data type}) podemos interrogar al arreglo para saber exactamente de qué estamos hablando.

\begin{nota}
    \textit{Recordatorio}: Si usted combina tipos de datos, \textit{NumPy} no se va a quejar, pero tomará decisiones por usted. Transformará todo al tipo de dato más "complejo" o flexible que encuentre. Si mete un \textit{float} en una lista de enteros, todos serán \textit{floats}. Si mete un \textit{string}, prepárese, todos los números se convertirán en texto y perderán sus propiedades matemáticas.
\end{nota}

El código no es muy complejo, vease:
\begin{pycode}{.dtype para tipo de dato}
import numpy as np 

# Creamos un array simple
array = np.array([1,2,3])
print(f'Tipo de dato: 'array.dtype) # devuelve: int
    
\end{pycode}

Es vital consultar \textit{.dtype} antes de realizar operaciones matemáticas complejas; no querrá intentar calcular una raíz cuadrada sobre un \textit{array} que \textit{NumPy} decidió convertir en texto porque usted olvidó una comilla por ahí.



\subsection{Conversión de tipos \textit{.astype()}}

Para finalizar este bloque, \textit{NumPy} nos trae una función que nos permite forzar el tipo de dato de un \textit{array} ya existente: \textit{.astype()}.

Suponga que está realizando una investigación y se encuentra con millones de datos de coma flotante (\textit{floats}). Usted determina que no tiene sentido cargar la memoria de su computadora con tanta precisión y cree que el error de conversión de decimal a entero no es significativo para su análisis. Pues bien, aquí es donde entra \textit{.astype()} para salvarle las papas.

También, siempre que sea posible, usted puede convertir texto a números, siempre y cuando ese texto sea realmente un número disfrazado de \textit{string}.

Vea el código en acción:

\begin{pycode}{Uso de .astype()}
import numpy as np

# Creamos un array de decimales (floats)
precios = np.array([10.5, 20.9, 30.2])
print(precios.dtype) # Resultado: float64

# Queremos solo la parte entera
precios_enteros = precios.astype(int)
print(precios_enteros) # Resultado: [10, 20, 30]
# (Ojo: no redondea, corta el decimal)

# Convertir a texto
precios_texto = precios.astype(str) 
print(precios_texto) # Resultado: ['10.5', '20.9', '30.2']

# Convertir de texto a float
precios_flotantes = precios_texto.astype(float)
print(precios_flotantes)

\end{pycode}

\begin{nota}
    \textbf{Advertencia:} Tenga mucho cuidado con el "hachazo". Al convertir de \textit{float} a \textit{int}, el valor 20.9 se convierte en 20, no en 21. Si su jefe le pregunta por qué faltan decimales en el presupuesto, no me eche la culpa a mí, yo le advertí que \textit{NumPy} es eficiente, pero no necesariamente piadoso con el redondeo.
\end{nota}




\subsection{Indexación y slicing}

Ahora que ya sabemos crear, deformar e interrogar a nuestros arreglos, llega el momento de aprender a operarlos.

La Indexación y el Slicing (o rebanado, para los que prefieren el castellano) son las herramientas que nos permiten extraer elementos específicos de un array. 

\begin{ejemplo}
Imagine que su matriz es una pizza gigante: la indexación es elegir un pepperoni exacto, y el slicing es llevarse una porción completa (o media pizza, si tiene hambre de datos).
\end{ejemplo}

\subsubsection{Indexación}

La indexación es el sistema de direcciones que utiliza \textit{NumPy}. Si un \textit{array} es un estante lleno de cajas, la indexación es la etiqueta que tiene cada una para que sea fácil encontrarla, sacarla y usar lo que hay dentro.

La indexación puede ser vectoria, matricial o tensorial;

Una indexación vectorial (recuerde la fila de su secundaria), es como señalar a alguien de la fila. Para señalar a alguien solo se necesita saber su posición numérica o índice. 

La indexcación matricial es cuando tenemos una tabla, recuerde el ejemplo de la reina en el tablero de agedrez, si usted la quisiera señalar deberia señalar la coordenada \(ij\), es decir la posición de la fila y columna (no recuerdo exactamente la posición, pero seria por ejemplo \(5C\)).

Por ultimo la indexación tensorial es un poco más compleja, imagine un cubo de Rubik; es como señalar un color específico en una de las caras internas del cubo. Aquí ya necesita tres coordenadas (o más, si es que usted disfruta del sufrimiento matemático).

Para utilizar la indexación, simplemente colocamos los corchetes inmediatamente después del nombre del \textit{array}: \texttt{array[indice]}.

Veamoslo actuar en código:

\begin{pycode}{Indexación en NumPy}
import numpy as np

# --- Caso 1D --- #
edades = np.array([20, 25, 30, 35])
print(edades[0])  # Resultado: 20
print(edades[-1]) # Resultado: 35 (el último)

# --- Caso 2D --- #
# Tabla de 2 filas y 3 columnas
ventas = np.array([
    [100, 120, 150], # Fila 0
    [200, 210, 250]  # Fila 1
    ])
dato = ventas[1, 0] # ventas del primer mes(columna), segundo año(fila)
print(dato) # Resultado: 200
\end{pycode}

Indexar no solo sirve para leer datos como si fuera un chisme; también sirve para escribir y corregir la realidad de su \textit{array}.

\begin{pycode}{Cambiando valores}
import numpy as np

precios = np.array([10.5, 20.0, 35.0])
# Queremos cambiar el segundo precio por 22 (ahora es 20)
precios[1] = 22.0

print(precios) # [10.5, 22.0, 35.0]
\end{pycode}

Si aún no termina de visualizarlo, considere este ejemplo culinario: 
\begin{ejemplo}
    \label{ejemplo_chef_lasana}
    El chef armo una lasaña de 9 porciones perfectas (3 filas y 3 columnas). Cada pedazo tiene un ingrediente secreto: uno tiene un borde quemado, otro mucha salsa, otro está fria, uno tiene demasiado queso, etc (Este chef no salio muy bueno ¡Se parece a mí!). Usted quiere elegir la de borde quemadito, pues no le parece nada cómico comer una porción con muchisima salsa, una porción fria o una que solo sabe a queso. Ademas le gusta esa textura crocante que deja el quemado. Pues usted cuando elije la porción esta haciendo una indexación (ademas de la ponderación) por lo tanto eligiria en el array lasaña el indice [i, j] donde se encuentra la lasañá con borde quemadito.
\end{ejemplo}

\emph{Indexación booleana}

La indexación no termina con las coordenadas. El creador de \textit{NumPy} sabía perfectamente que ni usted ni yo vamos a saber siempre en qué índice exacto se encuentra la porción de borde quemadito del ejemplo \ref{ejemplo_chef_lasana}. Por lo tanto, diseñó un "colador inteligente".

Se llama \textit{booleana} porque solo entiende dos estados: \textit{True} o \textit{False} (\textit{verdadero} o \textit{falso}, si es que nunca se dignó a abrir un libro de inglés).

La indexación booleana funciona en tres pasos:
\begin{enumerate}
\item \textbf{La pregunta:} Le preguntamos al \textit{array}: "¿Dónde la temperatura de la lasaña superó los 100 grados?".
\item \textbf{La máscara:} \textit{NumPy} analiza cada elemento y crea un nuevo arreglo de verdaderos y falsos (\textit{máscara}), marcando con \textit{True} solo los lugares que cumplen la condición.
\item \textbf{El filtrado:} Le pedimos a \textit{NumPy} que nos entregue el contenido original solo de los lugares donde la respuesta fue \textit{True}.
\end{enumerate}

Mas allá de su elección de la porcion de lasañá, en un análisis numérico serio (no más importante que elegir bien la porción de lasaña) esto nos ayuda a filtrar datos según lo que nos interesa: apartar valores atípicos (\textit{outliers}), segmentar clientes por edad o encontrar acciones que superaron cierto precio.

En código:

\begin{pycode}{Indexación Booleana paso a paso}
import numpy as np

edades = np.array([15, 25, 12, 40, 18])

# 1. La pregunta
# ¿Quiénes son mayores de edad (>= 18)?
# 2. La mascara o la busqueda de verdaderos y falsos
mascara = edades >= 18
print(mascara) # Resultado: [False,  True, False,  True,  True]
# Los pasos 1 y dos estan relacionados

#2. filtrar el array 
mayores = edades[mascara]
print(mayores) # Resultado: [25, 40, 18]

# --- Forma rápida (Todo en una línea) --- #

resultado_directo = edades[edades >= 18]

\end{pycode}

La indexación booleana también puede ser condicional. Si usted quiere filtros más complejos, de esos que requieren un "esto Y aquello" o "esto O lo otro", podemos usar operadores lógicos.

Ahora, Dios sabrá por qué \textit{NumPy} decidió no usar las palabras \textit{and}, \textit{or} y \textit{not} de \textit{Python} convencional (realmente es por una cuestión de eficiencia a nivel de bits), pero nos obliga a usar estos símbolos especiales:

\begin{itemize}
\item \textit{AND} (Y): Se usa el símbolo ``\&'' (ampersand).
\item \textit{OR} (O): Se usa el símbolo ``|'' (la barra vertical o \textit{pipe}).
\item \textit{NOT} (NO): Se usa el símbolo ``\textasciitilde'' (la virgulilla de la "ñ").
\end{itemize}

Veamos esto en código:

\begin{pycode}{Filtros complejos}

precios = np.array([10, 50, 100, 150, 200])
# Precios entre 60 y 180
# Cada condición debe ir entre paréntesis ( )
filtro = (precios > 60) & (precios < 180)

print(precios[filtro]) # Resultado: [100, 150]

\end{pycode}



\subsubsection{Slicing}
Si la indexación era señalar la porción de lasaña con el dedo, el \textit{slicing} (o rebanado) es meter el cuchillo y llevarse una sección entera utilizando el operador de los dos puntos (:).

El slicing consiste de la siguiente sintaxis o estructura básica:
\begin{verbatim}
    array([Inicio:Fin:Paso])
\end{verbatim}
\begin{nota}
    El fin al igual que en \textit{np.arange} no considera el valor final, es decir nunca lo toca (recuerde el intervalo semiabierto).
\end{nota}

El slicing vectorial es el más intuitivo. Imagine un \textit{array} de 6 elementos \((\left[0,1,2,3,4,5\right])\). Usted puede realizar los siguientes cortes:
\begin{itemize}
    \item array[1:4]: Te da los elementos en las posiciones 1, 2 y 3. (El 4 queda fuera).
    \item array[:3]: Los primeros tres (del inicio hasta el índice 2).
    \item array[3:]: Desde el índice 3 hasta el final.
    \item array[::2]: Todo el array, pero saltando de dos en dos (posiciones 0, 2, 4).
\end{itemize}

En el caso del slicing matricial, tenemos filas y columnas, por lo tanto usariamos comas para separar los cortes de cada dimensión:
\begin{verbatim}
    array([corte_filas,corte_columnas: ...])
\end{verbatim}
Se pueden realizar por ejemplo estos slicing:
\begin{itemize}
    \item Seleccionar una fila entera: array[1, :] (Fila índice 1, todas las columnas).
    \item Seleccionar una columna entera: array[:, 0] (Todas las filas, columna índice 0).
    \item Seleccionar una sub-tabla: array[0:2, 1:3] (Filas 0 y 1; columnas 1 y 2).
\end{itemize}

No nos vamos a meter a tensores, ya que nadie quiere perder la cabeza (ese nadie soy yo).

Veamos el slicing en acción en codigo:
\begin{pycode}{Slicing en NumPy}

import numpy as np

# Creamos una matriz de 4x4 (puedes imaginar meses vs productos)
datos = np.arange(16).reshape(4, 4)
"""
[[ 0,  1,  2,  3],
[ 4,  5,  6,  7],
[ 8,  9, 10, 11],
[12, 13, 14, 15]]
"""

# 1. Extraemos la segunda fila completa
fila_2 = datos[1, :] # [4, 5, 6, 7]

# 2. Extraer la última columna completa
ultima_col = datos[:, -1] # [3, 7, 11, 15]

#3. Extraer el cuadrado central (2x2)
centro = datos[1:3, 1:3]
"""
[[ 5,  6],
[ 9, 10]]
"""
\end{pycode}

\begin{nota}
    \textit{Atención}: El \textit{slicing} en \textit{NumPy} no crea un \textit{array} nuevo, crea una "vista" del original. Esto significa que si usted modifica la rebanada, \textit{también modifica el array original}. Si quiere un objeto independiente para hacer sus experimentos sin arruinar la fuente, use \textit{.copy()}.
\end{nota}
\section{Operaciones matemáticas}

Además de las operaciones vectoriales básicas que ya vimos, existen las funciones universales, mejor conocidas como \textit{ufuncs}. Lo más importante que debe saber es que estas funciones se aplican de manera vectorizada.

Si usted quisiera sacar la raíz cuadrada de 100 datos a mano, tendría que sacar la calculadora y empezar a pelearse con cada número. Pero como usted no quiere entregar su análisis un semestre más tarde, utiliza \textit{NumPy}, que aplica la operación a todo el vector, matriz o tensor elemento por elemento sin perder ni un ápice de eficiencia. ¡Si no fuera así, usaríamos la librería \textit{math} y nos iríamos a dormir una siesta mientras la computadora procesa!.

\subsection{Raiz cuadrada y exponencial}

Estas son las funciones "base" para una infinidad de cálculos estadísticos, financieros y científicos.
\begin{itemize}
    \item Raiz cuadrada: \textit{np.sqrt}, que calcúla \(\sqrt{x}\) para cada elemento.
    \item Exponencial: \textit{np.exp}: calcula \(e^x\) (donde \(e\) es el numero de Euler). Esta función es fundamental para calcular intereses compuestos, crecimientos poblacionales o para cuando su código decide crecer exponencialmente en errores. Podría dedicar toda una sección a explicar la belleza de\(e^x\), pero lamentablemente este no es un libro de álgebra y ambos tenemos cosas mejores que hacer.
\end{itemize}

Veamos estas operaciones en acción en su hábitat natural:
\begin{pycode}{Raiz y Exponencial}

import numpy as np

arr = np.array([1, 4, 9, 16])

# Raíz cuadrada
raices = np.sqrt(arr) # [1., 2., 3., 4.]

# Exponencial (e^x)
exp_valores = np.exp(np.array([0, 1, 2])) # [1., 2.718..., 7.389...]

\end{pycode}

Observe que los resultados son siempre \textit{floats} (decimales). NumPy sabe que las raíces y las exponenciales rara vez son exactas, así que se prepara para la imprecisión de la vida real desde el principio.



\subsection{Logaritmo Natural}

Si en matemáticas \(e^x\) nos dice cuánto creció una población en un tiempo determinado, el logaritmo natural \((ln(x))\) nos dice cuánto tiempo tuvo que pasar para alcanzar ese crecimiento.

Básicamente, el logaritmo natural es como un domador de monstruos numéricos. Se usa principalmente para:
\begin{itemize}
\item \textit{Normalizar escalas:} Si tiene datos que van de 1 a 1,000,000 (como seguidores en TikTok o la deuda externa), graficarlos es un caos visual. El logaritmo "aplasta" los números gigantes y "estira" los pequeños, haciendo que todo quepa elegantemente en una gráfica.
\item \textit{Linealizar relaciones:} Convierte curvas exponenciales complejas en líneas rectas, lo cual es oro puro para entrenar algoritmos de Inteligencia Artificial sin que su computadora explote en el intento.
\item \textit{Interés compuesto:} Fundamental para calcular cuánto tiempo tardará en ser millonario si decide invertir sus ahorros en lugar de gastárselos en café.
\end{itemize}
    

Imagine que se generó una colonia de bacterias que crece de forma masiva \((e^x)\) y usted quiere saber qué tan rápido fue ese crecimiento. Utilizando \textit{np.log()}, podemos estimar esos tiempos de la siguiente manera:

\begin{pycode}{Logaritmo natural}
import numpy as np

# Cantidad de bacterias en diferentes muestras
bacterias = np.array([1, 10, 100, 1000])

# Aplicamos el logaritmo para "suavizar" los datos
tiempos_estimados = np.log(bacterias)

print(tiempos_estimados)
# Resultado: [0. , 2.30258509, 4.60517019, 6.90775528]
\end{pycode}

Observe cómo pasamos de una escala que saltaba de 10 en 10 (multiplicando) a una escala lineal que avanza de forma constante. Eso es \textit{NumPy} poniendo orden en el caos.

\begin{nota}
    \textbf{Advertencia:} El logaritmo natural es una función muy útil, pero tiene un temperamento difícil: no acepta números negativos ni el cero. Si intenta calcular \textit{np.log(0)}, \textit{NumPy} le devolverá un \(-\infty\) (\textit{infinito negativo}) y un aviso de que algo está muy mal en su vida (o al menos en sus datos).
\end{nota}


\subsection{Funciones trigonométricas}

Imagine que tiene un asador de pollos giratorio. Las funciones trigonométricas no son más que las coordenadas que nos dicen dónde está el pollo en cada momento mientras da vueltas en un círculo unitario.

Imagine un círculo cuyo radio es exactamente 1. El pollo empieza a la derecha y gira en sentido contrario a las agujas del reloj. El ángulo (\(\theta\)) es qué tanto ha girado el asador.

\begin{tikzpicture}[>=stealth, scale=3]

    % Dibujar los ejes X e Y
    \draw[<->, thick] (-1.2, 0) -- (1.2, 0) node[right] {\textbf{x}};
    \draw[<->, thick] (0, -1.2) -- (0, 1.2) node[above] {\textbf{y}};

    % Dibujar el círculo unitario
    \draw[thick] (0,0) circle (1);

    % Definir variables para el ángulo
    \def\angle{45} % Puedes cambiar este valor para mover el punto
    \def\radius{1}

    % Dibujar el radio (línea desde el origen al punto)
    \draw[thick] (0,0) -- (\angle:\radius) node[midway, above left] {1};

    % Dibujar las líneas discontinuas (proyecciones)
    \draw[dashed] (\angle:\radius) -- ({\radius*cos(\angle)}, 0) node[below] {$\cos \theta$};
    \draw[dashed] (\angle:\radius) -- (0, {\radius*sin(\angle)}) node[left] {$\sin \theta$};

    % Dibujar el arco del ángulo theta
    \draw (0.2,0) arc (0:\angle:0.2);
    \node at (22.5:0.3) {$\theta$};

    % Dibujar los puntos negros en las intersecciones de los ejes
    \filldraw (1,0) circle (0.5pt) node[above right] {$(1, 0)$};
    \filldraw (-1,0) circle (0.5pt) node[above left] {$(-1, 0)$};
    \filldraw (0,1) circle (0.5pt) node[above right] {$(0, 1)$};
    \filldraw (0,-1) circle (0.5pt) node[below right] {$(0, -1)$};
    
    % Dibujar el punto sobre el círculo y su etiqueta
    \filldraw (\angle:\radius) circle (0.7pt) node[right, xshift=2pt] {$(\cos \theta, \sin \theta)$};

\end{tikzpicture}

Hay tres funciones principales:
\begin{itemize}
    \item Seno \(\sin\): Es la altura del pollo. Cuando el ángulo es 0°, la altura es 0. Cuando llega arriba (90°), la altura es 1.
    \item Coseno \(\cos \): Es la distancia horizontal desde el centro. Cuando el ángulo es 0°, el pollo está todo a la derecha (1). A los 90°, está justo en el centro horizontal (0).
    \item Tangente \(\tan\): Es la relación entre la altura y la distancia. Imagina la sombra que proyecta el pollo en una pared vertical lejana.
\end{itemize}

Si usted no es ingeniero ni arquitecto, podría pensar que esto no tiene utilidad en el análisis de datos. Se equivoca: la trigonometría está en todas partes, desde el procesamiento de audio y música, hasta el análisis de ciclos eléctricos, señales de radio y física de videojuegos.

\begin{nota}
    En \textit{NumPy}, las funciones esperan que les hables en Radianes, no en grados. \(180°= \pi\) radianes (aprox. 3.14159).
\end{nota}

Veamos cómo interrogar al pollo en código:

\begin{pycode}{Funciones trigonométricas}
import numpy as np

# Queremos saber la altura del pollo a los 90 grados
angulo_grados = 90
angulo_rad = np.radians(angulo_grados) # Convertimos a radianes

altura = np.sin(angulo_rad)
distancia = np.cos(angulo_rad)

print(f"A los 90°, la altura (Seno) es: {altura}") # Debería ser 1.0
print(f"A los 90°, la distancia (Coseno) es: {distancia}") # Debería ser casi 0
\end{pycode}

Como habrá notado, el coseno de 90° no dio exactamente 0, sino un número absurdamente pequeño \((6.12\times10^{-17})\). Esto sucede porque las computadoras tienen limitaciones de precisión decimal. No se asuste; para nosotros, eso es cero.

Existen miles de funciones como \textit{np.radians} o \textit{np.degrees} que no vamos a abarcar aquí. Si pretendiera explicar cada función de la librería, este apunte tendría más hojas que la documentación oficial y usted se quedaría sin espacio en su disco duro. Cuando necesite algo muy específico, no dude en consultar la documentación de \textit{NumPy}.


\section{Manipulación de arrays y filtrado}

Además del \textit{.reshape} visto en la sección de propiedades y la indexación y \textit{slicing}, existen otras formas de manipular y filtrar datos para que confiesen lo que necesitamos.

\subsection{Manipulación de arrays}

\subsubsection{Concatenación y apilamiento}

Si usted, como todo ser normal, viene del mundo de las hojas de cálculo, la concatenación de \textit{arrays} es simplemente la forma técnica de decir "copiar y pegar rangos de celdas" para armar una tabla más grande.

\paragraph{Concatenación (\textit{np.concatenate})}

Es la forma más general de unir dos o más arreglos. Esta función pega una lista de \textit{arrays} uno tras otro a lo largo de un eje específico.

Sin embargo, el autor sabe lo difícil que puede ser visualizar ejes en la cabeza sin que le empiece a doler la nuca, así que lo vamos a explicar en ``modo Umpa Lumpa''.

\begin{ejemplo}
Imagine que se encuentra en la fábrica de chocolates de Willy Wonka (esperemos que Paramount no me cobre derechos de autor). Usted y su pequeño compañero de diminuto tamaño están fabricando dos tipos de dulces: los rojos y los azules (dos arreglos vectoriales).
$$
\begin{aligned}
\text{\textit{Caramelos rojos}} &= [R, R, R] \
\text{\textit{Caramelos azules}} &= [A, A]
\end{aligned}
$$
Wonka grita: ¡Concatenen! Lo que usted hace es colocar los caramelos azules justo donde terminan los rojos, dándonos una sola fila:
\[
\left[R, R, R, A, A\right]
\]
\end{ejemplo}

Pero, mi pequeño compañero Umpa Lumpa, a veces la cosa se complica cuando dejamos los caramelos y pasamos a las bandejas (matrices) de \textit{brownies}.

Si utilizamos \textit{np.concatenate()} en matrices, debemos definir los \textit{axis} (ejes). El \textit{axis=0} hace que coloque una bandeja encima de la otra. Básicamente, la pila crece hacia arriba.

Pero tenga cuidado: apilar \textit{brownies} no es una actividad fácil cuando hablamos de millones de bandejas. Todas deben tener exactamente la misma dimensión en el ancho. No podemos poner una bandeja de cuatro \textit{brownies} sobre una de dos, porque se nos caería la torre y Wonka nos despediría sin indemnización.

\begin{nota}
    El \textit{axis=0} afecta a las filas. Si concatenamos dos matrices de \(2 \times 3\) usando \textit{axis=0}, obtendremos una de \(4\times 3\). El ancho se mantiene igual, pero la estructura ahora es más alta.
\end{nota}

``¿Qué pasa si en vez de ordenar por filas ordenamos por columnas?''. Pues lo único que debe hacer es utilizar \textit{axis=1}. Con esto le decimos a la máquina que pegue una bandeja al lado de la otra. Aquí también hay una condición: a Wonka (\textit{NumPy}) no le gusta que las bandejas tengan distinto largo (filas). Si intenta pegar una bandeja de 3 filas al lado de una de 2, la línea no encajaría y terminaría en la calle.

Véalo en el código de la fábrica:

\begin{pycode}{Concatenación}
import numpy as np

# Inventario de dulces
bandeja_1 = np.array([[1, 1], [1, 1]]) # 2x2
bandeja_2 = np.array([[2, 2], [2, 2]]) # 2x2

# Unir verticalmente (Hacer la torre más alta)
torre = np.concatenate((bandeja_1, bandeja_2), axis=0) # Resultado: Una matriz de 4x2

# Unir horizontalmente (Hacer la fila más larga)
fila_larga = np.concatenate((bandeja_1, bandeja_2), axis=1) # Resultado: Una matriz de 2x4
\end{pycode}

\begin{nota}
    \textit{Truco de sintaxis:} Fíjese que los arreglos que quiere unir van dentro de un par de paréntesis extra \texttt{((a, b))}. Esto es porque \textit{np.concatenate} espera una "secuencia" (como una tupla o lista) de \textit{arrays}, no los \textit{arrays} sueltos.
\end{nota}

Ahora usted entiende que es cocatenar una matriz y un vector (sin la parte matemática). 

Pero ¿Para qué sirve esto fuera de la fábrica de chocolates? Usted, que probablemente no tiene un contrato vigente con el señor Wonka, se preguntará: ``¿Y para qué quiero saber esto?''. La respuesta es que, en el análisis de datos, la concatenación es su mejor amiga cuando los registros vienen por separado.

Imagine estas situaciones cotidianas para un analista: 
\begin{itemize}
    \item Suponga que su sistema genera un registro de ventas diferente cada día. Si usted quiere realizar un análisis semanal, no va a mirar siete archivos por separado como un monje medieval; usted va a concatenar esos siete registros para tener un reporte semanal sólido en un solo \textit{array}.
    \item \textit{Fusión de fuentes:} Si tiene dos bases de datos que hablan del mismo tema (por ejemplo, encuestas de dos sucursales distintas), al concatenarlas fortalece su muestra. Una muestra más grande significa una estadística más precisa y una media de datos que no miente tanto.
\end{itemize}

En resumen, concatenamos para unificar. Un dato suelto es una anécdota, mil datos son una estadística.


\paragraph{Apilamiento}

Seamos sinceros, estar peleandose con los ejes (\textit{axis=0, axis=1}) es una actividad de alto riesgo para su salud mental. \textit{NumPy} sabe que usted es un ser humano (o algo parecido) y le ofrece ``atajos'' semánticos que hacen lo mnismo que concatenate, pero con nombres que hasta un niño de cinco años entenderia.

Si \textit{concatenate} es la herramienta multiuso, \textit{vstack} y \textit{hstack} son herramientas especializadas. Sus nombres vienen del inglés \textit{vertical stack} (apilamiento vertical) y \textit{horizontal stack} (apilamiento horizontal).

\begin{itemize}
\item \textit{np.vstack((a, b)):} Es el equivalente a \textit{concatenate} con \textit{axis=0}. Apila los arreglos verticalmente, uno sobre otro, como si estuviera armando una torre de panqueques.
\item \textit{np.hstack((a, b)):} Es el equivalente a \textit{concatenate} con \textit{axis=1}. Apila los arreglos horizontalmente, uno al lado del otro, como si estuviera formando una fila en el supermercado.
\end{itemize}

Fíjese qué limpio queda el código cuando dejamos de hablar de "ejes" y empezamos a hablar de "direcciones":

\begin{pycode}{Atajos de apilamiento}
import numpy as np

a = np.array([1, 1, 1])
b = np.array([2, 2, 2])

# Apilamos hacia arriba (Vertical)
vertical = np.vstack((a, b)) # Resultado: una matriz de 2x3

# Apilamos hacia el lado (Horizontal)
horizontal = np.hstack((a, b)) # Resultado: un vector largo de 6 elementos

print("Vertical:\n", vertical)
print("Horizontal:", horizontal)
\end{pycode}

\begin{nota}
    \textit{Un detalle de vital importancia:} Al igual que con \textit{concatenate}, estas funciones esperan recibir una tupla (es decir, doble paréntesis \texttt{((a, b))}). Si olvida el paréntesis interno, NumPy pensará que le está pasando argumentos extra y le lanzará un error que le arruinará la tarde.
\end{nota}




\subsubsection{División de arrays (\textit{split})}

Si concatenar y apilar era el arte de construir bloques, \textit{np.split()} es la herramienta para desarmarlos. Esta función corta un \textit{array} en partes exactamente iguales y genera una lista de sub-arreglos más pequeños.

A diferencia del \textit{slicing}, que es como sacar una rebanada de pan, \textit{split} está diseñado para repartir todo el bloque de un solo comando.

\begin{nota}
    \textit{La regla de la división exacta:} \textit{np.split} es extremadamente estricto. Si usted le pide dividir un \textit{array} en tres partes, el tamaño total del arreglo debe ser divisible por 3. Si la división no da un número entero, \textit{NumPy} entrará en pánico y lanzará un error porque no sabe qué hacer con el elemento que sobra. 
\end{nota}

\textit{np.split()} tiene el siguiente modo de operar:
\begin{verbatim}
    np.split(array, cantidad de divisones)
\end{verbatim}

Veamos cómo se descuartiza un vector en el código:

\begin{pycode}{Uso de np.split}
import numpy as np

# Creamos un vector de 6 elementos
datos = np.array([10, 20, 30, 40, 50, 60])

# Dividir en 3 partes iguales
pedazos = np.split(datos, 3)

print(pedazos[0]) # [10, 20]
print(pedazos[1]) # [30, 40]
print(pedazos[2]) # [50, 60]
\end{pycode}

En el caso de las matrices, volvemos a encontrarnos con nuestros viejos amigos, los \textit{axis} (ejes), para elegir la dirección del "hachazo":
\begin{itemize}
\item \textit{axis=0}: Es un corte horizontal. Divide la matriz a través de sus filas (obtenemos trozos de arriba y abajo).
\item \textit{axis=1}: Es un corte vertical. Divide la matriz a través de sus columnas (obtenemos trozos de izquierda y derecha).
\end{itemize}

Vea el siguiente ejemplo:
\begin{pycode}{Split en Matrices}
import numpy as np

matriz = np.arange(16).reshape(4, 4)

# Dividir la matriz a la mitad (horizontalmente)
arriba, abajo = np.split(matriz, 2, axis=0) # Nos quedan dos matrices de 2x4

# Dividir la matriz a la mitad (verticalmente)
izquierda, derecha = np.split(matriz, 2, axis=1) # Nos quedan dos matrices de 4x2

print(arriba, abajo)
print()
print(izquierda, derecha)

\end{pycode}

¿Para qué sirve esto en la vida real? Usted usará \textit{np.split()} principalmente para dos cosas:
\begin{enumerate}
\item \textit{Separación de variables:} Si tiene una tabla donde las primeras columnas son características (como altura y peso) y la última es el precio, puede usar \textit{split} para separarlas y procesarlas por rutas distintas.
\item \textit{Procesamiento por lotes (\textit{batches}):} Si tiene un millón de registros y su computadora empieza a echar humo, puede dividirlos en 10 partes iguales y procesar una por vez para darle un respiro al procesador.
\end{enumerate}

\begin{nota}
Para facilitarnos la vida una vez más, \textit{NumPy} también incluyó las funciones \textit{np.vsplit} (vertical) y \textit{np.hsplit} (horizontal), que funcionan exactamente igual que los atajos de apilamiento, evitándole tener que escribir el \textit{axis} manualmente.
\end{nota}



\subsubsection{Transposición y permutación de ejes}

\paragraph{Transposición}

Transponer una matriz consiste en el arte de intercambiar sus filas por sus columnas. Si usted tenía una fila, ahora tiene una columna; si tenía una columna, ahora es una fila. Es, básicamente, rotar su tabla de datos 90 grados y darle la vuelta.

Matemáticamente, si un elemento estaba en la posición \(i,j\) donde \(i\) es una fila y \(j\) es una columna, la matriz transpuesta ahora sera \(j,i\), Esto se visualiza de la siguiente manera:
\[
A = \begin{bmatrix}
    1 & 2 \\
    3 & 4 \\
    5 & 6
\end{bmatrix}
\quad
\text{Transpuesta}
\quad
A^T= \begin{bmatrix}
    1 & 3 & 5 \\
    2 & 4 & 6
\end{bmatrix}
\]

¿De qué sirve transponer una matriz en la vida real?
\begin{itemize}
\item \textit{Reorientar datos:} Si sus sensores registraron datos en columnas pero su algoritmo los espera en filas, la transposición lo arregla en un milisegundo.
\item \textit{Álgebra Lineal:} Es fundamental para cumplir las condiciones de la multiplicación de matrices (que veremos si el destino lo permite).
\item \textit{Correlaciones:} Permite comparar cómo se relacionan diferentes variables entre sí al alinear correctamente los ejes.
\end{itemize}

Para transponer, \textit{NumPy} ofrece la función \texttt{np.transpose(array)}, pero como usted es un programador que valora su tiempo, usará la propiedad \texttt{.T}. Es corta, es elegante y le ahorra preciosos segundos de vida que puede gastar en ver videos de gatitos o tomar café.

Véalo en el código:
\begin{pycode}{Transposicion con .T}
import numpy as np

# Creamos una matriz de 2 filas y 3 columnas
matriz = np.array([
[1, 2, 3],
[4, 5, 6]
])

print(f"Forma original: {matriz.shape}") # (2, 3)

# Transponemos la matriz
matriz_t = matriz.T

print(matriz_t)
"""
Resultado:
[[1, 4],
[2, 5],
[3, 6]]
"""
print(f"Nueva forma: {matriz_t.shape}") # (3, 2)
\end{pycode}

\begin{nota}
    Si intentas transponer un vector de una sola dimensión (un array simple como \texttt{[1, 2, 3]}), verás que \textit{no pasa nada}. Esto es porque NumPy considera que un vector 1D no tiene filas ni columnas que intercambiar. Para "parar" un vector, primero deberías convertirlo en una matriz de \(1 \times n\) usando \texttt{.reshape()}.
\end{nota}



\paragraph{Permutación de ejes}

Si la transposición (\texttt{.T}) es como dar vuelta toda la media, \texttt{np.swapaxes()} es un movimiento mucho más preciso: intercambia exactamente dos ejes entre sí, dejando todos los demás intactos en su posición original. Su estructura consiste en:
\begin{verbatim}
    np.swapaxes(array, eje_1, eje_2)
\end{verbatim}


Esta función es la favorita de quienes trabajan con datos multidimensionales, como imágenes o tensores complejos, donde a veces solo necesita rotar dos dimensiones específicas (por ejemplo, intercambiar el ancho por el alto) sin alterar la profundidad o el número de canales de color.

En una matriz simple de dos dimensiones, usar \textit{swapaxes} es exactamente lo mismo que usar \textit{.T}. La verdadera diferencia —y donde usted agradecerá su existencia— aparece cuando trabajamos con tres o más dimensiones.

Vea este ejemplo:

\begin{pycode}{transpose vs swapaxes}
import numpy as np

# Creamos un "cubo" de datos (2 matrices de 3x4)
cubo = np.zeros((2, 3, 4)) # np.shape: (0: profundidad, 1: filas, 2: columnas)

# 1. Usando transpose para invertir TODO
cubo_t = np.transpose(cubo) # La forma pasará a ser (4, 3, 2)

# 2. Usando swapaxes para cambiar solo FILAS por COLUMNAS
cubo_swapped = np.swapaxes(cubo, 1, 2) # La forma pasará a ser (2, 4, 3)
# Solo tocamos el eje 1 y el eje 2. La profundidad (eje 0) queda igual.

print(f"Original: {cubo.shape}")
print(f"Transpose: {cubo_t.shape}")
print(f"Swapaxes: {cubo_swapped.shape}")
\end{pycode}

Como puede observar, \texttt{swapaxes} es mucho más respetuoso con la estructura que usted no quiere tocar. Mientras que \texttt{transpose} lo pone todo patas arriba.




\subsection{filtrado de datos}

Además de la indexación y la indexación booleana existe el filtrado condicional con \texttt{np.where()}. Esta función es, en esencia, un bloque \textit{if-else} convertido en un superatleta vectorizado.

Su estructura consiste en:
\begin{verbatim}
    np.where(condición, valor_verdadero, valor_falso).
\end{verbatim}

Suponga que quiere evaluar una lista de alumnos para separar el trigo de la paja (aprobados de desaprobados). En lugar de ir uno por uno preguntando la nota, usted le da la orden a NumPy y él se encarga del resto:
\begin{pycode}{Uso basico de np.where}
import numpy as np

notas = np.array([3, 8, 5, 2, 10])

# Si la nota es >= 4, ponemos 1 (Aprobado), sino 0 (Desaprobado)
resultados = np.where(notas >= 4, 1, 0)

print(resultados) #Resultado: [0, 1, 1, 0, 1]
\end{pycode}

Existen dos razones fundamentales por las que usted, si aprecia su tiempo, usará \texttt{np.where} en lugar de un \texttt{if} tradicional de \textit{Python}:
\begin{itemize}
\item \textit{Vectorización (Velocidad):} Como mencionamos antes, NumPy procesa todos los elementos al mismo tiempo. Si tuviera un millón de filas, \textit{np.where} tardaría milisegundos, mientras que un bucle \textit{for} con un \textit{if} le daría tiempo de ir a prepararse un café (y quizás hasta de cultivarlo).
\item \textit{Limpieza de Datos:} Es la herramienta ideal para corregir errores absurdos en sus \textit{datasets}. Por ejemplo: "Si el precio es negativo, póngale 0; de lo contrario, deje el precio original".
\end{itemize}

Vea cómo limpiar su honor (y sus datos) con un solo comando:
\begin{pycode}{Limpieza de datos con np.where}

precios = np.array([100, -50, 200, -10])

# Si el precio es menor a 0, lo corrijo a 0. Si no, mantengo el precio original.
precios_limpios = np.where(precios < 0, 0, precios)

print(precios_limpios) # Resultado: [100, 0, 200, 0]
\end{pycode}

Si usted le pasa a \texttt{np.where} solo la condición (sin los valores de verdadero/falso), la función cambia de personalidad: en lugar de devolver valores transformados, le devolverá los índices exactos donde se cumple lo que usted pidió.

\begin{pycode}{np.where para encontrar indices}
import numpy as np

edades = np.array([10, 25, 18, 40])

indices = np.where(edades >= 18) # Resultado: (array([1, 2, 3]),)
print(indices)
\end{pycode}

\begin{nota}
Fíjese que el resultado es una tupla que contiene un \textit{array}. Esto parece innecesariamente complejo, pero es así porque está preparado para matrices de muchas dimensiones donde usted recibiría una lista de índices por cada eje. Para nosotros, por ahora, basta con saber que ahí están nuestras coordenadas.
\end{nota}
\section{Estadística}

\textit{NumPy} incluye funciones estadísticas extremadamente rápidas, ya que operan sobre \textit{arrays} completos o sobre subconjuntos definidos por ejes, lo que permite analizar millones de datos sin escribir un solo bucle manual.

\subsection{Mínimo y máximo (\textit{np.min} y \textit{np.max})}

Para su felicidad y la mía, estas funciones no requieren un doctorado para entenderlas: simplemente nos entregan el valor más \textit{bajo} y el más \textit{alto} de un conjunto de datos.

\begin{pycode}{Mínimo y Máximo}
import numpy as np

# Creamos el array
arr = np.array([5, 12, 7, 20, 15])

print("Mínimo:", np.min(arr)) # 5
print("Máximo:", np.max(arr)) # 20


# También existen los métodos .min() y .max() que funcionan igual:
print("Mínimo (método):", arr.min())
print("Máximo (método):", arr.max())
\end{pycode}



\subsection{Media y mediana (\textit{np.mean} y \textit{np.median})}

Aquí es donde la estadística empieza a ponerse interesante.
\begin{itemize}
\item \textit{Media (Promedio):} Es la suma de todos los valores dividida por la cantidad de elementos. Es útil cuando sus datos son "estables" y no tienen locuras en los extremos.
\item \textit{Mediana:} Es el valor que queda justo en el centro cuando ordena los datos de menor a mayor. Es el mejor amigo del analista cuando hay valores atípicos (\textit{outliers}) que intentan mentirle.
\end{itemize}
\begin{ejemplo}
        Imagine que cinco amigos ganan 1.000 y el sexto es un millonario que gana 1.000.000. La media dirá que el promedio de sueldos es altísimo (una mentira), mientras que la mediana dirá que el sueldo central es 1.000 (la realidad de la mayoría).
\end{ejemplo}

Vea un ejemplo:
\begin{pycode}{Calculo de media y mediana}
import numpy as np

# Creamos el array
arr = np.array([1, 2, 3, 4, 5, 100])

# Dejamos que NumPy calcule la media
print("Media:", np.mean(arr)) # 19.166...

# Calculo de la mediana
print("Mediana:", np.median(arr)) # 3.5
\end{pycode}

\subsection{Desviación estándar (\textit{np.std})}

La desviación estándar mide qué tan "desparramados" están los datos respecto al promedio. Es como cuando usted dice que llegará a las 8:00, "más o menos 10 minutos". Esos 10 minutos son su desviación.

Basicamente nos dice qué tan "lejos" están los datos respecto al promedio.
\begin{itemize}
\item \textit{Baja desviación:} Los datos son estables y predecibles (están cerca de la media).
\item \textit{Alta desviación:} Hay mucha variabilidad o riesgo (los datos están por todos lados).
\end{itemize}

Vea un ejemplo:
\begin{pycode}{Desviación estandar}
import numpy as np

# Creamos el array
arr = np.array([10, 12, 14, 16, 18])

# Obtenemos la media 
print("Media:", np.mean(arr)) # 14

# Vemos la dispersión de los datos
print("Desviación estándar:", np.std(arr)) # 2.828...
\end{pycode}



\subsection{El dominio de los ejes (\textit{axis})}

El eje o axis es la dirección en la que \textit{NumPy} realiza la operación.

\begin{itemize}
\item \textit{axis=0 (Hacia abajo):} La operación atraviesa las filas. El resultado es un resumen de cada columna.
\item \textit{axis=1 (Hacia los lados):} La operación atraviesa las columnas. El resultado es un resumen de cada fila.
\end{itemize}

Matematicamente una matriz se forma por \textit{filas} y \textit{columnas} (\(i,j\)). El indice 0 corresponde a las filas, el indice 1 corresponde a las columnas.

Vea un ejemplode la utilización de \texttt{axis} para obtener la media y la mediana.
\begin{pycode}{Utilziación de axis}
import numpy as np

# Filas: Estudiantes | Columnas: Materias 
notas = np.array([
    [8, 10], # Estudiante A
    [7, 9],  # Estudiante B
    [9, 6]   # Estudiante C
])

# 1. ¿Cuál es el promedio de cada MATERIA? (Bajar por las filas)
promedio_materias = np.mean(notas, axis=0) 
# Resultado: [8.0, 8.33] -> [Promedio Mates, Promedio Arte]

# 2. ¿Cuál es el promedio de cada ESTUDIANTE? (Moverse por las columnas)
promedio_estudiantes = np.mean(notas, axis=1)
# Resultado: [9.0, 8.0, 7.5] -> [Promedio A, Promedio B, Promedio C]
\end{pycode}
\section{Álgebra lineal con \textit{NumPy}}

Llegamos a la sección donde NumPy se pone el traje de gala y saca la calculadora científica: el Álgebra Lineal. Si usted piensa que esto es solo para matemáticos con mucho tiempo libre, se equivoca; esta es la base de cómo una IA reconoce su cara en una foto o cómo Google decide qué mostrarle primero.

En \textit{NumPy}, casi todo lo que sea "pesado" matemáticamente vive dentro del submódulo \texttt{np.linalg}.

\subsection{Suma y producto (\textit{np.sum} y \textit{np.prod})}

Estas funciones básicas lo que hacen es ``colapsar'' su arreglo en un único valor. \texttt{np.sum()} suma todos los elementos de su \texttt{array}, mientras que \textit{np.prod()} multiplica todos los elementos entre sí.
\begin{pycode}
import numpy as np

# Creamos un array 
arr = np.array([1, 2, 3, 4])

# Calculamos la suma de los elementos
print("Suma:", np.sum(arr)) # 10

# Calculamos el producto de los elementos
print("Producto:", np.prod(arr)) # 24
\end{pycode}

\begin{nota}
    \textit{np.sum()} es extremadamente útil en la indexación booleana. Como los \textit{True} valen 1 y los \textit{False} valen 0, si usted hace \texttt{np.sum(edades >= 18)}, obtendrá automáticamente la cantidad de personas mayores de edad en su arreglo.
\end{nota}

\subsection{Producto punto (\textit{np.dot()})}

Matemáticamente, el producto punto (o producto escalar) de dos vectores consiste en multiplicar los elementos que ocupan la \textit{misma posición} y luego sumar todos esos resultados.

Para realizar esta operación, hay una regla de etiqueta inquebrantable: ambos vectores deben tener el mismo número de elementos. Si uno tiene tres datos y el otro tiene cuatro, \textit{NumPy} se negará a trabajar.

Si tenemos el vector \(A = \begin{bmatrix}
    a_1 & a_2
\end{bmatrix}\) y el vector \(B = \begin{bmatrix}
    b_1 & b_2
\end{bmatrix}\) el producto punto de \(A, B\) (\(A\centerdot B\)) es:
\[
A \centerdot B = (a_1 \times b_1) + (a_2 \times b_2)
\]

En analisis de datos esto es un poco más sencillo, ya que consta solo del comando:
\begin{verbatim}
    np.dot(vector1, vector2)
\end{verbatim}

Usted me preguntara ``¿Para qué sirve esto en la vida real?'' Aunque no lo crea, esta es una herramienta tan poderosa que, si aprende a usarla, será su pan de cada día:

En comercio o finanzas sirve para calcular ingreso total. Imagine que tiene tres porductos, y que vendio las siguentes cantidades \([10, 5, 20]\)  aun precio unitario de \([100, 250, 50]\), el producto punto nos da el ingreso total en un solo paso:
\[
(10 \times 100) + (50 \times 250) + (20 \times 50) = 1000 + 1250+ 1000 = 3250
\]

Es muy utilizado en redes sociales. Los sistemas de recomendación crean un vector con sus gustos y un vector con el contenido disponible. Si el producto punto entre sus gustos y el vector son altos, significa que son parecidos, por lo tanto lo va a recomendar

Hay dos formas de calcular el producto punto:
\begin{enumerate}
    \item Con \texttt{np.dot(vector1, vector2)}
    \item Con \texttt{@}, \texttt{vector1 @ vector2}
\end{enumerate}

Vea el siguiente ejemplo:
\begin{pycode}{Producto Punto en NumPy}
import numpy as np

# Creación de vectores 
cantidades = np.array([10, 5, 20])
precios = np.array([100, 250, 50])

# Operación 1: Funcion explicita
total_1 = np.dot(cantidades, precios)

# Operación 2: Operador @ 
total_2 = cantidades @ precios

print(f"El ingreso total es: {total_2}") # 3250
\end{pycode}

\subsubsection{Producto de matrices}

El producto punto también funciona en matrices. Para multiplicar dos matrices, no pueden tener cualquier forma. Existe una condición matemática estricta: El número de columnas de la primera matriz debe ser igual al número de filas de la segunda.

A diferencia del producto punto de vectores, en el caso de la matriz no es una multiplicación elemento a elemento. El proceso es: Fila por Columna.

Imagine que usted tiene una matriz de Recetas \((2 platos \times 3 ingredientes)\) y una matriz de Costos \((3 ingredientes \times 1 precio)\).
\begin{itemize}
\item Plato A necesita: [2 huevos,1 harina,3 leches].
\item Plato B necesita: [1 huevo,2 harinas,1 leche].
\end{itemize}
El producto punto \texttt{(@)} multiplicará cada fila de la "Receta" por la columna de "Costos" y sumará los resultados. El resultado final será una matriz de 2 platos por 1 costo total.

Véalo en el código, donde la magia (y la matemática) sucede:

\begin{pycode}{Multiplicación de Matrices}
import numpy as np

# Matriz A: 2 filas, 3 columnas (2x3)
# [Huevos, Harina, Leche] para dos platos
recetas = np.array([
[2, 1, 3],
[1, 2, 1]
])

# Matriz B: 3 filas, 1 columna (3x1)
# Precio de cada ingrediente
precios = np.array([
[50],
[80],
[120]
])

# El producto punto matricial (@)
costo_total = recetas @ precios

print("Costo total por plato:\n", costo_total)

# Resultado:
# [[540]   <- Costo Plato A
# [330]]  <- Costo Plato B
\end{pycode}

\begin{nota}
    En la multiplicación de matrices, el orden \textit{sí} altera el producto. \texttt{A@B} no es lo mismo que \texttt{B@A}. De hecho, es muy probable que si intenta hacer \texttt{B@A}, las dimensiones ni siquiera coincidan/
\end{nota}

Además de las ya vistas, hay otra forma de multiplicar matrices, que funciona igual que producto punto, este metdodo es \texttt{np.matmul(A,B)}.

\paragraph{Multiplicación elemento a elemento}
El producto de \textit{Hadamard} es, básicamente, una multiplicación "espejo". Esta operación toma el número de una posición en la Matriz A y lo multiplica exclusivamente por el número que ocupa la misma posición exacta en la Matriz B.

Para que esto funcione sin que \textit{NumPy} le grite, las matrices deben ser gemelas: deben tener exactamente la misma forma(por ejemplo, ambos de \(2 \times 3\))

Matemáticamente, si el elemento \(A[0,0] \times B[0,0] = C[0,0]\) el segundo elemento de \(C[0,1]\) es \(A[0,1] \times B[0,1]\).

\begin{ejemplo}
    Suponga que usted tiene una matriz de Precios de tres productos en dos sucursales distintas, y otra matriz con las Cantidades vendidas en esos mismos lugares. Al hacer \texttt{Precios * Cantidades}, usted obtiene una nueva matriz que le indica la \textit{ganancia total por producto y sucursal}.
\end{ejemplo}

En \textit{NumPy}, esta operación es la que viene "por defecto" cuando usa el asterisco:
\begin{verbatim}
    A * B
\end{verbatim}

Véalo en el código para no confundirlo con el producto punto:

\begin{pycode}{Hadamard vs Producto Punto}
import numpy as np

# Creación de matriz A y B
A = np.array([[1, 2], [3, 4]])
B = np.array([[10, 20], [30, 40]])

# 1. Producto de Hadamard (Elemento a elemento)
hadamard = A * B # Resultado: [[110, 220], [330, 440]]
print("Hadamard (*):\n", hadamard)


# 2. Producto Punto (Matricial - Filas x Columnas)
matricial = A @ B # Resultado: [[110 + 230, 120 + 240], ...]
print("Matricial (@):\n", matricial)
\end{pycode}

\paragraph{Hadamard vs. Producto Punto}

Para que no se le haga un nudo en las neuronas, imagine que sus matrices son grupos de personas. El producto de Hadamard es como si cada persona se mirara en un espejo; el Producto Punto es como si todo el grupo hiciera una coreografía coordinada.

Use el Producto de Hadamard (*) cuando:
\begin{itemize}
\item \textit{Aplicación de Impuestos o Descuentos:} Si tiene una matriz de precios y una matriz de porcentajes de descuento, los multiplica uno a uno para obtener el precio final de cada producto específico.
\item \textit{Filtros de Imágenes:} En computación visual, si quiere oscurecer una parte de una foto, multiplica los píxeles por una "máscara" de números menores a 1.
\item \textit{Cambio de Unidades:} Convertir una matriz de metros a centímetros multiplicando todo por un arreglo de 100.
\end{itemize}

Use el Producto Punto (@) cuando:
\begin{itemize}
\item \textit{Cálculo de Ingresos Totales:} Como vimos con el chef, cuando quiere cruzar cantidades vendidas con una lista de precios para obtener un gran total.
\item \textit{Redes Neuronales (IA):} Las neuronas artificiales reciben señales y las multiplican por "pesos" para decidir si se activan o no. Esta combinación de señales es puro producto punto.
\item \textit{Sistemas de Recomendación:} Para medir qué tan parecidos son sus gustos con los atributos de una película en Netflix.
\end{itemize}




\subsection{Matriz identidad e inversa}

\subsubsection{Matriz identidad}

Si el número 1 es el elemento neutro de la aritmética básica \((5\times1=5)\), la Matriz Identidad es el elemento neutro del álgebra de matrices.

Esta tiene mucha utilidad en álgebra lineal, ya que no existe la división de matrices. Para dividir una matriz por otra la multiplicamos por la inversa. La matriz identidad es la pieza clave para calcularla: Una matriz \(A\) multiplicada por su inversa \(A^{-1}\) da como resultado la identidad.

Véalo en el código para refrescar la memoria:

\begin{pycode}{La Matriz Identidad en acción}
import numpy as np

# Creamos una identidad de 3x3
I = np.eye(3)

#Creamos una matriz cualquiera
A = np.array([
[10, 20, 30],
[40, 50, 60],
[70, 80, 90]
])

# Multiplicación matricial por la identidad (Use @, no *)
resultado = A @ I
# El resultado es exactamente A. La identidad no cambió nada.
print(np.array_equal(A, resultado)) # Resultado: True
\end{pycode}


\subsubsection{Matriz inversa}

Como mencioné anteriormente —y si es que todavía sigue prestando atención—, en el mundo de las matrices no existe el símbolo de dividir. Por lo tanto, si usted necesita realizar el equivalente a una división, lo que realmente necesita es la inversa.

La matriz inversa (\(A^{-1}\)) es aquella que al multiplicarse por la original da como resultado la matriz identidad.
\[
A \times A^{-1} = I
\]

En términos aritméticos simples: si tenemos el número \(5\), su inverso es \(1/5\), porque \(5 \times 1/5 = 1\).

\begin{ejemplo}
    Imagine que tiene una máquina procesadora de frutas. Usted introduce una banana perfecta (\(X\)). La máquina está programada para pelarla y rebanarla (operación \(A\)), lo que nos entrega un plato de rodajas de banana (\(B\)).

    Si usted se arrepiente del proceso y quiere que esas rodajas vuelvan a ser una banana entera, necesitaría una "máquina inversa" ($A^{-1}$). Si tal maravilla de la ingeniería existiera, al meter las rodajas ($B$), la máquina debería devolverle la banana original ($X$).
\end{ejemplo}

No se emocione demasiado, no todas las matrices tienen inversa. Si una matriz tiene un determinante igual a 0, se le llama "singular". Es el equivalente matemático a que su banana haya pasado por una licuadora: la información se perdió para siempre y \textit{NumPy} le lanzará un error (\textit{LinAlgError}) si intenta invertirla.

Para encontrar esta "máquina del tiempo" en \textit{NumPy}, utilizamos el submódulo de álgebra lineal:
\begin{verbatim}
    np.linalg.inv(matriz) 
\end{verbatim}

Vea esto en codigo:
\begin{pycode}{Calcular Matriz Inversa}
import numpy as np

# Definimos una matriz cuadrada
A = np.array([[1, 2], [3, 4]])

# Calculamos la inversa
A_inv = np.linalg.inv(A)

print("Matriz Inversa:\n", A_inv) # Comprobamos que A @ A_inv da la Identidad

identidad = np.round(A @ A_inv) # Usamos round para limpiar decimales minimos de precision

print("\nResultado (Identidad):\n", identidad) # Devuelve la matriz identidad
\end{pycode}





\subsection{Determinantes y autovalores}

\subsubsection{Determinante (\textit{np.linalg.det})}

El determinante es un valor escalar (un solo número) que se calcula únicamente para matrices cuadradas (aquellas con el mismo número de filas que de columnas). Si su matriz es rectangular, ni lo intente: el determinante no existe allí.

Matemáticamente, para una matriz de \(2 \times 2\), el cálculo es un simple cruce de caminos: \textit{multiplicamos} las \textit{diagonales} y las \textit{restamos}.
\[A = \begin{bmatrix}
    a & b \\
    c & d
\end{bmatrix} = 
det(A) = (a \times d) - (b \times c)
\]

¿Para qué sirve este número? La utilidad del determinante va mucho más allá de ser un ejercicio de álgebra, es el guardián de la Inversa:
\begin{itemize}
\item \textit{Si el determinante es 0:} La matriz es \textit{singular}. No tiene inversa.
\item \textbf{Si el determinante es distinto de 0:} La matriz es \textit{regular}. Tiene inversa
\end{itemize}

Calcular determinantes de matrices de \(4 \times 4\) o más es una verdadera pesadilla que ha hecho llorar a generaciones. Por suerte, \textit{NumPy} lo resuelve con un solo comando:
\begin{verbatim}
    np.linalg.det(matriz)
\end{verbatim}

Veamos esto en acción:
\begin{pycode}{Calculo del Determinante}
import numpy as np

# Matriz A (Tiene inversa porque su det no sera 0)
A = np.array([[4, 2], [1, 3]])

# Calculamos el determinante de A
det_A = np.linalg.det(A) 
print(f"Determinante de A: {det_A}") # Calculo manual: (4*3) - (2*1) = 10 

# Matriz B (Lineas dependientes, det sera 0)
B = np.array([[1, 2], [2, 4]])

# Calculamos el determinante de B
det_B = np.linalg.det(B)
print(f"Determinante de B: {det_B}") # Resultado: 0.0 (No tiene inversa)
\end{pycode}

\subsubsection{Autovalores y autovectores (\textit{np.linalg.eig})}

Imagine que una matriz es una fuerza de la naturaleza que transforma el espacio: puede estirar, rotar o deformar un plano de datos. En medio de ese caos, existen elementos que mantienen la compostura:

\begin{itemize}
\item \textbf{Autovector:} Es una dirección especial que, tras aplicar la transformación, no cambia su orientación. Puede estirarse o encogerse, pero sigue apuntando exactamente hacia el mismo lugar. Es el "eje" o la columna vertebral de la transformación.
\item \textbf{Autovalor:} Es el número que nos indica cuánto se estiró o encogió ese autovector. Si el autovalor es 2, el vector duplicó su tamaño. Si es 0.5, se redujo a la mitad.
\end{itemize}

En \textit{NumPy}, utilizamos el comando \texttt{np.linalg.eig()}, donde \textit{eig} proviene del alemán \textit{eigen}, que significa "propio" o "característico".

Esto en codigo se ve como:
\begin{pycode}{Calcular Autovalores y Autovectores}
import numpy as np

# Definimos una matriz de transformacion
A = np.array([[3, 1], [0, 2]])

# Obtenemos ambos resultados
valores, vectores = np.linalg.eig(A)

print("Autovalores (Eigenvalues):", valores) # Resultado: [3., 2.] -> Indica los factores de estiramiento

print("\nAutovectores (Eigenvectors):\n", vectores) # Cada columna es un vector que mantiene su direccion
\end{pycode}


Este concepto es bastante dificil, por lo que vamos a ver con manzanitas.

\paragraph{Autovector (la dirección)}

Imagione que pone una manzana en una prensa. La prensa aplasta desde arriba. Si mira la manzana, su piel se estira hacia los lados y se deforma. Sin embargo hay una linea imaginaria (vector) que atraviesa la manzana justo por el centro, de arriba a abajo. Mientras todo lo demas se desparrama, la linea vertical sigue apuntando exactamente en la misma dirección. No se desplazo hacia ningun lado. Es decir, es la dirección que, aunque la matriz (la prensa) la estruje, sigue apuntando hacia la misma dirección. Es el "eje" de la transformación.

\paragraph{Autovalor}

Una vez identificado el autovector, miramos qué le pasó a su tamaño. Si la prensa redujo la altura de la manzana a la mitad, el autovalor es \(0.5\). Si la máquina fuera un estirador y duplicara su altura, el autovalor sería \(2\). Si no cambia el tamaño, el autovalor es \(1\).

\paragraph{Función}

Imagine un camión con 1,000 manzanas y demasiada información: peso, color, diámetro, dulzor, firmeza, etc. (demasiadas dimensiones).

\begin{enumerate}
\item \textit{Encontrar el Autovector principal:} El algoritmo busca la "dirección" en la que los datos varían más. Quizás descubre que el Peso y el Diámetro se mueven siempre juntos.
\item \textit{Usar el Autovalor:} Si el autovalor de esa dirección es muy grande, significa que esa combinación (Peso + Diámetro) explica casi todo sobre sus manzanas.
\end{enumerate}

Resultado: En lugar de analizar 5 variables, ahora solo mira esa "Súper Variable". Ha simplificado su problema (reducción de dimensionalidad) sin perder la esencia de la información.


\section{Números aleatorios y semillas}

Bienvenido a la última parte de este texto. Siéntase orgulloso: después de haber dominado temas que le darían pesadillas a un matemático medieval, llegamos a algo más sencillo.

\subsection{Números aleatorios}

Aunque Python trae su propia librería random, la versión de NumPy es mucho más potente porque está diseñada para trabajar con arreglos gigantes de un solo golpe y con una eficiencia envidiable.
\begin{verbatim}
    np.random()
\end{verbatim}


\subsubsection{Generar numeros aleatorios entre 0 y 1}

Esta función es el estándar para generar "probabilidades". Devuelve valores en un intervalo semiabierto \([0,1)\), lo que significa que puede salir 0, pero nunca llegará a ser exactamente 1.
\begin{verbatim}
    np.random.rand()
\end{verbatim}

Vealo con un ejemplo
\begin{pycode}{Número entre 0 y 1}
    import numpy as np

    print("Número aleatorio [0,1):", np.random.rand())
\end{pycode}

\subsubsection{Array de números aleatorios}

Si usted quisiera generar un array con números aleatorios utilice la función:
\begin{verbatim}
    np.random.rand()
\end{verbatim}

En código:
\begin{pycode}{Array aleatorio}
    import numpy as np

    print("Array 1D:", np.random.rand(5)) # 5 números aleatorios
    print("Matriz 2D:\n", np.random.rand(2, 3)) # Matriz 2x3
\end{pycode}

\subsubsection{Enteros aleatorios}

Esta función es la que usted usará cuando quiera simular el lanzamiento de un dado o elegir un ID de usuario al azar. Su estructura es muy clara:
\begin{verbatim}
    np.random.randint(inicio, fin (no lo incluye), cantidad)
\end{verbatim}

En código:
\begin{pycode}{Enteros aleatorios}
    import numpy as np
    
    print("Enteros aleatorios:", np.random.randint(1, 10, size=5))
    # 5 números entre 1 y 9
\end{pycode}


\subsection{Fijar semillas \textit{np.random.seed()}}

Las computadoras, por muy inteligentes que parezcan, son máquinas deterministas: no son realmente capaces de crear azar de la nada. Lo que hacen es ejecutar una fórmula matemática tan compleja que genera una lista de números que parecen aleatorios. A esto lo llamamos números \textit{pseudoaleatorios}.

La semilla (\textit{seed}) es el número inicial que se le entrega a esa fórmula. Si usted usa la misma semilla, la fórmula siempre arrancará en el mismo lugar y, por lo tanto, toda la secuencia de números siguientes será exactamente igual.

Pero ¿Cómo funciona? Imagine que tiene un libro gigante con mil millones de números escritos.
\begin{itemize}
\item \textbf{Sin semilla:} \textit{Python} abre el libro en una página cualquiera (basándose usualmente en el reloj interno de su computadora). Cada vez que corre el programa, el libro se abre en un lugar distinto.
\item \textbf{Con semilla \textit{np.random.seed(42)}:} Usted le ordena a \textit{Python}: "Abra siempre el libro en la página 42". No importa cuántas veces reinicie su computadora, la secuencia de números empezará siempre en el mismo renglón.
\end{itemize}

Véalo usted mismo: si corre este código, obtendrá los mismos decimales que yo:
\begin{pycode}{Efecto de la semilla}
import numpy as np

# Ponemos una semilla cualquiera (ej: 42)
np.random.seed(42)
print("Primer intento:", np.random.rand(3)) # [0.37454012, 0.95071431, 0.73199394]

#Si volvemos a poner la misma semilla...
np.random.seed(42)
print("Segundo intento:", np.random.rand(3)) # [0.37454012, 0.95071431, 0.73199394]
# Los resultados seran IDENTICOS.
\end{pycode}

¿Para qué nos sirve? En el análisis de datos, la aleatoriedad es necesaria pero peligrosa. Aquí es donde la semilla nos salva:
\begin{enumerate}
\item \textit{Reproducibilidad (La clave de la ciencia):} Si usted descubre una tendencia increíble en una simulación financiera, necesita que su jefe o su profesor puedan correr su código y ver exactamente lo mismo que vio usted. Sin la semilla, ellos obtendrían otros números y pensarían que usted se inventó las conclusiones.
\item \textit{Caza de errores (\textit{Debugging}):} Si su programa falla solo una de cada diez veces que genera datos al azar, arreglarlo es una pesadilla porque el error "desaparece" al reiniciar. Al fijar la semilla, el error se vuelve estable y usted puede diseccionarlo con calma.
\item \textit{Comparación justa:} Si quiere saber qué estrategia de inversión es mejor, debe probar ambas bajo las mismas crisis y bonanzas. La semilla le permite crear un escenario "al azar" idéntico para todos sus modelos.
\end{enumerate}



\subsection{Distribuciones comunes estadísticas}

En el análisis de datos, las distribuciones son modelos matemáticos que describen con qué frecuencia ocurre cada valor en un conjunto de datos. No todos los datos se comportan igual: algunos prefieren amontonarse en el centro y otros se reparten democráticamente.

\subsubsection{Distribucion Uniforme}

Es aquella en la que todos los resultados posibles tienen exactamente la misma probabilidad de ocurrir. Visualmente, en un histograma, se ve como un rectángulo perfectamente plano. No hay "campana" ni picos; es una meseta de probabilidad donde el azar no tiene favoritos. Vease la figura \ref{distribución_uniforme_numpy}

Dependiendo de si necesita números enteros o decimales, \textit{NumPy} le ofrece tres herramientas principales:

\begin{enumerate}
\item \textit{Uniforme Continua (\texttt{np.random.rand} o \texttt{np.random.uniform}):} Genera números decimales. Se usa cuando el dato puede tener infinitos valores intermedios, como el tiempo exacto de espera de un buque en el puerto o la tasa de interés diaria de un bono.
\item \textit{Uniforme Discreta (\texttt{np.random.randint}):} Genera números enteros. Se usa cuando el dato es contable, como elegir un contenedor al azar de un lote de 100 o simular el número de una factura.
\end{enumerate}


\begin{figure}[!ht]
\centering
\begin{tikzpicture}
  \begin{axis}[
    axis lines = middle,
    xlabel = {$x$},
    ylabel = {$f(x)$},
    ymin = 0, ymax = 0.6,
    xmin = -1, xmax = 6,
    xtick = {1, 4},
    xticklabels = {$a$, $b$},
    ytick = {0.333},
    yticklabels = {$\frac{1}{b-a}$},
    every axis x label/.style={at={(current axis.right of origin)},anchor=north west},
    every axis y label/.style={at={(current axis.above origin)},anchor=south east},
    axis line style={->},
  ]

  % Área sombreada bajo la curva
  \addplot [fill=blue!20, draw=none, domain=1:4] {1/3} \closedcycle;

  % Líneas de la distribución uniforme
  \addplot [very thick, blue, domain=1:4] {1/3};
  
  % Líneas verticales discontinuas para los límites a y b
  \draw [dashed, gray] (axis cs:1,0) -- (axis cs:1,0.333);
  \draw [dashed, gray] (axis cs:4,0) -- (axis cs:4,0.333);

  % Líneas en cero para x < a y x > b
  \addplot [very thick, blue, domain=-0.5:1] {0};
  \addplot [very thick, blue, domain=4:5.5] {0};

  \end{axis}
\end{tikzpicture}
\caption{Representación de una distribución uniforme}
\label{distribución_uniforme_numpy}
\end{figure}

\paragraph{¿Para qué sirve en la vida real?}
\begin{itemize}
\item \textbf{Sorteos y Muestreos:} Si quiere auditar 10 documentos de un total de 500, usa esta distribución para asegurar que el documento 1 tenga la misma chance de ser revisado que el 500.
\item \textbf{Simulaciones Iniciales:} Cuando no tiene idea de cómo se comporta una variable (por ejemplo, el precio de un nuevo \textit{commodity}), empieza asumiendo que puede valer cualquier cosa entre un mínimo y un máximo con la misma probabilidad.
\item \textbf{Criptografía:} Las claves de seguridad se generan buscando una distribución uniforme perfecta para que un atacante no pueda predecir qué números salen más seguido.
\end{itemize}

Veamos una distribucion uniforme en código:
\begin{pycode}{Distribucion Uniforme}

import numpy as np

# 1. Decimales entre 0 y 1 (5 valores)
decimales = np.random.rand(5)

# 2. Decimales en un rango especifico (ej: entre 10 y 50)
precios_azar = np.random.uniform(10, 50, size=5)

# 3. Enteros (Simular el lanzamiento de un dado 10 veces)
dados = np.random.randint(1, 7, size=10) # El limite superior (7) no se incluye
\end{pycode}




\subsubsection{Distribución normal o Gausseana}

Aquí los datos tienen un favorito: el promedio. Es la forma que toman casi todos los fenómenos naturales y sociales cuando tienes una muestra lo suficientemente grande. En NumPy, la generamos principalmente con \texttt{np.random.randn()} o \texttt{np.random.normal()}.

Se la conoce como la Campana de Gauss. Su característica principal es la simetría. La mayoría de los datos se concentran en el centro (donde están la media, la mediana y la moda), a medida que te alejas del centro hacia la izquierda o derecha, la frecuencia de los datos cae rápidamente. Los valores "extremos" (muy altos o muy bajos) son muy poco probables. Vea la figura \ref{distribución_gausseana_numpy}.

\pgfmathdeclarefunction{gauss}{3}{%
  \pgfmathparse{1/(#3*sqrt(2*pi))*exp(-((#1-#2)^2)/(2*#3^2))}%
}
\begin{figure}[!ht]
    \centering
    \begin{tikzpicture}
  \begin{axis}[
    no markers, 
    domain=-4:4, 
    samples=100,
    axis lines=left, 
    xlabel=$x$, ylabel=$f(x)$,
    every axis y label/.style={at=(current axis.above origin),anchor=south},
    every axis x label/.style={at=(current axis.right of origin),anchor=west},
    height=5cm, width=10cm,
    xtick={-3, -2, -1, 0, 1, 2, 3},
    ytick=\empty,
    enlargelimits=false, 
    clip=false, 
    axis on top,
    grid = none
  ]

  % Área bajo la curva (68% - Una desviación estándar)
  \addplot [fill=blue!20, draw=none, domain=-1:1] {gauss(x, 0, 1)} \closedcycle;

  % Curva de la distribución normal
  \addplot [very thick, cyan!50!black] {gauss(x, 0, 1)};

  % Etiquetas de parámetros
  \draw [dashed] (axis cs:0,0) -- (axis cs:0,{gauss(0,0,1)}) 
    node[above] {$\mu$};
    
  % Indicación de Sigma
  \draw [latex-latex] (axis cs:0, 0.2) -- (axis cs:1, 0.2) 
    node[midway, above] {$\sigma$};

  \end{axis}
\end{tikzpicture}
\caption{Distribución normal o gausseana}
\label{distribución_gausseana_numpy}
\end{figure}

La distribución normal tiene dos numeros importantes:
\begin{itemize}
    \item La media(\(\mu\)): Indica donde esta el centro o promedio de la campana.
    \item La desviación estandar (\(\omega\)): Indica que tan ancha o flaca es la campana. Si la desviación es pequeña todos los datos se encuentran pegados al promedio, si es grande los datos estan más dispersos.
\end{itemize}

En análisis de datos, usamos la distribución normal porque nos permite predecir el futuro con márgenes de error:
\begin{itemize}
    \item Control de Calidad: Si una máquina llena botellas de 500ml, el llenado real seguirá una distribución normal. Si una botella sale con 450ml, sabemos que es un "error extraño" porque está muy lejos del centro de la campana.
    \item Finanzas y Riesgo: Los rendimientos de las acciones suelen modelarse así. La desviación estándar nos dice qué tan arriesgada es una inversión.
    \item Biometría y Sociología: La altura de las personas, los coeficientes intelectuales o los errores de medición en experimentos científicos.
\end{itemize}

Veamos como se realiza la campana de gauss en \textit{NumPy}
\begin{pycode}{Distribucion Normal}
import numpy as np

# 1. Normal Estandar (Media = 0, Desviacion = 1)
datos_estandar = np.random.randn(1000) # Generamos 1000 puntos para que se note la forma

# 2. Normal Personalizada (Media, Desviacion, Cantidad)
alturas = np.random.normal(175, 10, 1000) # Altura promedio 175cm con dispersion de 10cm

# 3. Ver el promedio real de los datos generados
print(f"Media real: {np.mean(alturas)}") # Deberia ser muy cercano a 175
\end{pycode}




\subsubsection{Distribución binomial}

Es la que usamos cuando los datos no son medidas (como la altura o el peso), sino conteos de éxitos o fracasos. Es la distribución del "Sí o No", del "Pasa o no Pasa". En NumPy, la generamos con \texttt{np.random.binomial()}.

Imagina que realizas un experimento (como lanzar una moneda) un número determinado de veces (\(n\)). En cada intento, solo hay dos resultados posibles: Éxito o Fracaso. La distribución binomial te dice cuál es la probabilidad de obtener exactamente \(k\) éxitos en esos n intentos, conociendo de antemano la probabilidad \(p\) de que un solo intento sea exitoso.


Para que NumPy calcule esta distribución, necesita tres parámetros:
\begin{itemize}
    \item \textit{n (n)}: El número de veces que repites el experimento (ej: lanzas la moneda 10 veces).
    \item \textit{p (prob)}: La probabilidad de éxito en un solo intento (ej: 0.5 para una moneda, o 0.05 para una pieza fallada).
    \item \textit{size}: Cuántas veces quieres repetir todo el conjunto de experimentos para ver los resultados.
\end{itemize}
    
¿Para qué sirve?  En comercio y control de procesos, es la herramienta fundamental para la toma de decisiones, ya que permite:
\begin{itemize}
    \item Control de Calidad: Si sabes que el \(2\%\) de las piezas de una fábrica suelen salir falladas, la distribución binomial te dice qué tan probable es que en un lote de 100 piezas encuentres más de 5 falladas (lo cual indicaría que algo anda mal en la máquina).
    \item Conversión de Ventas: Si un vendedor tiene una tasa de cierre del \(10\%\), ¿cuál es la probabilidad de que logre 3 ventas si llama a 20 clientes hoy?
    \item Logística: Si la probabilidad de que un camión se retrase en la frontera es del \(15\%\), ¿cuántos camiones de una flota de 10 llegarán tarde probablemente?
\end{itemize}
    

En \textit{NumPy }esto se ve así:
\begin{pycode}{Distribucion Binomial}
import numpy as np

# Ejemplo: Lanzamos una moneda 10 veces (n=10)
# La probabilidad de cara es 0.5 (p=0.5)
# Queremos simular este juego 1 vez
exitos = np.random.binomial(n=10, p=0.5, size=1) 
print(f"Caras obtenidas: {exitos}")

# Ejemplo de Negocios:
# Probabilidad de que un cliente pague a término: 0.8 (p)
# Tenemos 100 clientes (n)
# Simulamos 1000 meses para ver el comportamiento (size)
pagos_a_termino = np.random.binomial(n=100, p=0.8, size=1000)

print(f"Promedio de clientes que pagan a término: {np.mean(pagos_a_termino)}")
\end{pycode}

\end{document}
